\documentclass[a4paper,11pt,english]{article}

\usepackage[margin=2.5cm]{geometry}
\usepackage[utf8]{inputenc}
\usepackage[T1]{fontenc}
\usepackage[english]{babel}
\usepackage{amsmath}
\usepackage{amssymb}
\usepackage{amsthm}
\usepackage[colorlinks,urlcolor=blue,citecolor=blue,linkcolor=blue]{hyperref}

\newtheorem{theorem}{Theorem}[section]
\newtheorem{lemma}[theorem]{Lemma}
\newtheorem{corollary}[theorem]{Corollary}
\newtheorem{assumption}[theorem]{Assumption}

\newcommand{\E}{\mathbb E}
\newcommand{\Prob}{\mathbb P}
\newcommand{\Var}{\operatorname{Var}}
\newcommand{\Cov}{\operatorname{Cov}}
\newcommand{\norm}[1]{\left\lVert #1\right\rVert}
\newcommand{\R}{\mathbb R}

\setlength{\emergencystretch}{3em}

\title{Moment Approximation for Specification Search}
\author{Th\'eo Voldoire and Nic Fishman}
\date{}

\begin{document}

\maketitle

\begin{abstract}
Sparse specification search is a smooth search over Schur complements of
one empirical second-moment matrix.  Expanding this field before Gaussian
replacement separates its quadratic term into a covariance-matched bias,
a negligible diagonal fluctuation, and an off-diagonal field that is
linear in each incoming row.  This gives a covariance-matched Gaussian
approximation throughout $s\log(ep/s)=o(n)$ while retaining arbitrary
observational drift and signal.  Under sparse treatment alignment the
same field reduces to bilinear score products.  Their conditional
increments yield a covariance-aware forward search with an exact regret
certificate; a block-Kronecker Gaussian score law gives vanishing regret
without local isotropy, while local isotropy permits approximation by a
static top-$s$ sort.  The same
moment field gives a matched studentized approximation and unified upper
and lower bounds, with lower tightness governed by metric separation.
\end{abstract}

\section{The sparse moment field}

For independent and identically distributed observations
$i=1,\ldots,n$, let $x^{(i)}$ be a centered focal regressor,
$X^{(i)}\in\R^p$ centered controls, and $y^{(i)}$ a centered outcome.
Put
\begin{equation*}
        \beta_\emptyset^*=\frac{\E[x^{(i)}y^{(i)}]}{\E[(x^{(i)})^2]},
        \qquad
        r^{(i)}=y^{(i)}-\beta_\emptyset^*x^{(i)}.
\end{equation*}
Then $\E[x^{(i)}r^{(i)}]=0$.  Set
\begin{equation*}
        \E[(x^{(i)})^2]=\sigma_0^2,
        \qquad
        \E[(r^{(i)})^2]=\sigma_r^2,
        \qquad
        \E[X^{(i)}(X^{(i)})^\top]=C,
\end{equation*}
and assume $\sigma_0\sigma_r>0$.  Discard zero-variance controls.
Rescaling the remaining controls changes none of their selected spans, so
$C_{jj}=1$ throughout.

The augmented row and its second-moment matrix are
\begin{equation}\label{eq:moment-matrix}
        v^{(i)}=
        \begin{pmatrix}x^{(i)}\\ r^{(i)}\\ X^{(i)}\end{pmatrix},
        \qquad
        \Gamma=\E[v^{(i)}(v^{(i)})^\top]
        =\begin{pmatrix}
        \sigma_0^2&0&d^\top\\
        0&\sigma_r^2&h^\top\\
        d&h&C
        \end{pmatrix},
\end{equation}
where $d=\E[X^{(i)}x^{(i)}]$ and $h=\E[X^{(i)}r^{(i)}]$.  Let
\begin{equation}\label{eq:empirical-moment-matrix}
        \hat\Gamma
        =\frac1n\sum_{i=1}^n v^{(i)}(v^{(i)})^\top.
\end{equation}

Write $\mathbf x,\mathbf r,\mathbf y\in\R^n$ for the stacked observations
and let $\mathbf X\in\R^{n\times p}$ have rows $(X^{(i)})^\top$.
All sample projections below use these bold symbols; unbolded
$(x,r,y,X)$ denotes a generic observation.

For an empirical or population positive semidefinite control block,
every inverse and inverse square root below acts on its range;
equivalently, it is the Moore--Penrose inverse in the original
coordinates.  Perturbations of a singular population block are first
compressed to that fixed range, where the inverse is ordinary.  For a
symmetric matrix $A$, two indices $a,b\in\{x,r\}$, and a control set
$S$, define the residual moment
\begin{equation}\label{eq:residual-moment}
        A_{ab\mathbin{\cdot}S}
        =A_{ab}-A_{aS}A_{SS}^{-1}A_{Sb}.
\end{equation}
The coefficient contrast associated with $S$ is the function
\begin{equation}\label{eq:coefficient-functional}
        B_S(A)
        =\frac{A_{xr\mathbin{\cdot}S}}{A_{xx\mathbin{\cdot}S}}
        -\frac{A_{xr}}{A_{xx}},
        \qquad B_\emptyset(A)=0.
\end{equation}
On a sample for which a scalar denominator in
\eqref{eq:coefficient-functional} vanishes, set the corresponding empirical
contrast to zero.  This measurable convention only defines the field on the
exceptional event; Theorem~\ref{thm:critical-empirical-expansion} shows that
all searched denominators are positive with the stated probability.
This function carries both the observational population landscape and the
stochastic search around it.  It is invariant under replacing $r$ by
$r+c x$ for any scalar $c$, because both ratios in
\eqref{eq:coefficient-functional} then increase by $c$.  The empirical
field can therefore be computed from $(\mathbf x,\mathbf y,\mathbf X)$ even though the orthogonal
residual $r$ makes its population geometry clearer.

For $S\subseteq[p]$, let $M_S$ project orthogonally off the columns of
$\mathbf X_S$.  When the denominator is positive, write
\begin{equation*}
        \hat\beta_{0,S}
        =\frac{\mathbf x^\top M_S\mathbf y}
        {\mathbf x^\top M_S\mathbf x}.
\end{equation*}
Set $\hat\beta_{0,\emptyset}=0$ when $\mathbf x^\top\mathbf x=0$, and
set $\hat\beta_{0,S}=\hat\beta_{0,\emptyset}$ whenever
$\mathbf x^\top M_S\mathbf x=0$.
When $\Gamma_{xx\mathbin{\cdot}S}>0$, let $\beta_{0,S}^*$ denote the
coefficient on $x$ in the population linear projection of $y$ on
$(x,X_S)$; equivalently,
\begin{equation*}
        \beta_{0,S}^*
        =\beta_\emptyset^*
        +\frac{\Gamma_{xr\mathbin{\cdot}S}}
        {\Gamma_{xx\mathbin{\cdot}S}}.
\end{equation*}

\begin{lemma}[Exact moment representation]\label{lem:exact-moment}
For every support with $\mathbf x^\top M_S\mathbf x>0$,
\begin{equation}\label{eq:exact-moment-representation}
        \hat\beta_{0,S}-\hat\beta_{0,\emptyset}
        =B_S(\hat\Gamma).
\end{equation}
For every support with $\Gamma_{xx\mathbin{\cdot}S}>0$, the population
matrix satisfies
\begin{equation}\label{eq:population-landscape}
        B_S(\Gamma)
        =\beta_{0,S}^*-\beta_{0,\emptyset}^*
        =-\frac{d_S^\top C_{SS}^{-1}h_S}
        {\sigma_0^2-d_S^\top C_{SS}^{-1}d_S}.
\end{equation}
\end{lemma}

\begin{proof}
Since $\mathbf y=\beta_\emptyset^*\mathbf x+\mathbf r$, the coefficient on
$\mathbf x$ after controlling for $\mathbf X_S$ is $\beta_\emptyset^*$ plus
the coefficient on $\mathbf x$ in the regression of $\mathbf r$ on
$(\mathbf x,\mathbf X_S)$.  The partitioned-regression formula makes
the latter coefficient
$\hat\Gamma_{xr\mathbin{\cdot}S}/
\hat\Gamma_{xx\mathbin{\cdot}S}$.  Subtracting the corresponding
coefficient with $S=\emptyset$ proves
\eqref{eq:exact-moment-representation}.  Substitution of
\eqref{eq:moment-matrix} into \eqref{eq:coefficient-functional} proves
\eqref{eq:population-landscape}.
\end{proof}

The exact field depends only on selected column spans, so singular
population blocks require no additional assumption.  Indeed, write
$C_{JJ}=E_J\Lambda_JE_J^\top$ on its positive-eigenvalue range.  If
$a\in\ker(C_{JJ})$, then $\E(a^\top X_J^{(i)})^2=0$, hence
$a^\top X_J^{(i)}=0$ almost surely.  Thus
\begin{equation*}
        \operatorname{span}(\mathbf X_J)
        =\operatorname{span}(\mathbf X_JE_J\Lambda_J^{-1/2}),
        \qquad
        \E[\widetilde X_J^{(i)}(\widetilde X_J^{(i)})^\top]=I,
        \quad
        \widetilde X_J^{(i)}
        =\Lambda_J^{-1/2}E_J^\top X_J^{(i)}.
\end{equation*}
The vectors $d_J,h_J$ and both the empirical and matched Gaussian
perturbations lie in this same fixed range.  We therefore carry out every
inverse and derivative in these whitened coordinates and return to the
original coordinates with the Moore--Penrose inverse.

Let $\Pi_J=E_JE_J^\top$ be the projection onto the range of
$C_{JJ}$.  A symmetric perturbation $H$ is range-compatible if
\begin{equation}\label{eq:range-compatible}
        H_{Jx},H_{Jr}\in\operatorname{range}(C_{JJ}),
        \qquad H_{JJ}=\Pi_JH_{JJ}\Pi_J,
        \qquad |J|\le2s.
\end{equation}
Every empirical row-moment perturbation and every matched Gaussian
perturbation is range-compatible.  Indeed, every forbidden kernel
functional of an empirical row moment vanishes almost surely, so its
variance is zero; covariance matching forces the corresponding Gaussian
functional to vanish as well.  All derivatives below are taken on
this fixed-rank affine face.  When a function is evaluated on an ambient
matrix, its support-specific blocks are projected onto the face first.

The population residual treatment variance is
\begin{equation}\label{eq:overlap}
        \eta(S)
        =\Gamma_{xx\mathbin{\cdot}S}
        =\sigma_0^2-d_S^\top C_{SS}^{-1}d_S.
\end{equation}
\begin{assumption}[Sparse moment regime]\label{ass:sparse-moments}
For a fixed constant $K$, the augmented row satisfies
\begin{equation*}
        \norm{a^\top v^{(i)}}_{\psi_2}
        \le K(a^\top\Gamma a)^{1/2}
        \qquad\text{for every }a\in\R^{p+2},
\end{equation*}
and
\begin{equation*}
        \inf_{|S|\le s}\eta(S)\ge c_{\mathrm{ov}}\sigma_0^2
\end{equation*}
for a constant $c_{\mathrm{ov}}>0$.
\end{assumption}

Overlap is the domain condition of
\eqref{eq:coefficient-functional}; relative sub-Gaussianity supplies
uniform moment concentration on the compressed sparse blocks.
Under the linear specialization below, it follows from relative
sub-Gaussianity of the design together with independent noise that is
sub-Gaussian relative to its variance.

Write
\begin{equation}\label{eq:search-scale}
        L=\log(ep/s),
        \qquad m=sL,
        \qquad \delta=\delta_{s,n,p}=\sqrt{m/n}.
\end{equation}
The regime $m=o(n)$ makes the empirical sparse blocks a vanishing relative
perturbation of their population values.

At the critical scale $m\lesssim n$, a second-order empirical expansion
controls every support.  Under $m=o(n)$, the searched field admits a
covariance-matched Gaussian replacement.  Under sparse treatment
alignment, this Gaussian field loses its Gram component and reduces to
bilinear score products.

\section{Matched moment and score approximations}

\subsection{Matched moments}

The empirical matrix in \eqref{eq:empirical-moment-matrix} is an average
of row outer products.  Let $G$ be the centered symmetric Gaussian matrix
with matched covariance
\begin{equation}\label{eq:matched-moment-matrix}
        \Cov(G_{ab},G_{cd})
        =\Cov(v_a^{(i)}v_b^{(i)},v_c^{(i)}v_d^{(i)}),
        \qquad a,b,c,d\in\{x,r,1,\ldots,p\}.
\end{equation}
The covariance in \eqref{eq:matched-moment-matrix} retains the fourth
moments of the original rows.  Its second-order Taylor field therefore
retains the population landscape, the treatment and outcome scores, the
control Gram fluctuations, and their quadratic interaction.

For a symmetric perturbation $H$, define its relative sparse size by
\begin{equation}\label{eq:sparse-moment-size}
\begin{split}
        \lvert H\rvert_s
        ={}&\frac{|H_{xx}|}{\sigma_0^2}
        +\frac{|H_{xr}|}{\sigma_0\sigma_r}\ +
        \sup_{|J|\le2s}\bigg\{
        \frac{\norm{C_{JJ}^{-1/2}H_{Jx}}}{\sigma_0}
        +\frac{\norm{C_{JJ}^{-1/2}H_{Jr}}}{\sigma_r}\\
        &\hspace{43mm}
        +\norm{C_{JJ}^{-1/2}H_{JJ}C_{JJ}^{-1/2}}_{\mathrm{op}}
        \bigg\}.
\end{split}
\end{equation}

For $z\in\R^p$, define
\begin{equation*}
        \norm{z}_{C,k}
        =\sup_{|J|\le k}
        \left(z_J^\top C_{JJ}^{-1}z_J\right)^{1/2},
\end{equation*}
and put
\begin{equation}\label{eq:mean-sizes}
        \phi_k=\frac{\norm{d}_{C,k}^2}{\sigma_0^2},
        \qquad
        \widetilde\Psi_k=\norm{h}_{C,k},
        \qquad
        a_n=\frac1{\sigma_0}
        \{(\sigma_r\sqrt{\phi_s}+\widetilde\Psi_s)\delta
        +\sigma_r\delta^2\}.
\end{equation}
The first two terms in $a_n$ bound the treatment and outcome parts of the
linear field; the last is the size of the quadratic field when that
derivative degenerates.

For real random variables $A,B$, let
\begin{equation*}
        d_3(A,B)
        =\sup_f|\E f(A)-\E f(B)|,
\end{equation*}
where the supremum is over three times continuously differentiable $f$
with $\max_{0\le j\le3}\norm{f^{(j)}}_\infty\le1$.

Define the empirical and matched second-order fields
\begin{align}
        \mathcal T_{n,S}
        &=B_S(\Gamma)+DB_S(\Gamma)[\hat\Gamma-\Gamma]
        +\frac12D^2B_S(\Gamma)
        [\hat\Gamma-\Gamma,\hat\Gamma-\Gamma],
        \label{eq:empirical-second-order-field}\\
        \mathcal T^G_{n,S}
        &=B_S(\Gamma)+\frac1{\sqrt n}DB_S(\Gamma)[G]
        +\frac1{2n}D^2B_S(\Gamma)[G,G].
        \label{eq:gaussian-taylor-field}
\end{align}

\begin{theorem}[Critical empirical expansion]
\label{thm:critical-empirical-expansion}
Suppose Assumption~\ref{ass:sparse-moments} holds and
$\delta\le\delta_0$ for a sufficiently small constant.  With probability
at least $1-Ce^{-cm}$, every searched coefficient is defined and,
simultaneously over $|S|\le s$,
\begin{align}
\left|B_S(\hat\Gamma)-B_S(\Gamma)
-DB_S(\Gamma)[\hat\Gamma-\Gamma]\right|
&\le C\frac{\sigma_r}{\sigma_0}\delta^2,
\label{eq:critical-first-order-expansion}\\
\left|B_S(\hat\Gamma)-\mathcal T_{n,S}\right|
&\le C\frac{\sigma_r}{\sigma_0}\delta^3.
\label{eq:critical-second-order-expansion}
\end{align}
Consequently,
\begin{equation}\label{eq:critical-field-maximum}
\left|
\max_{|S|=s}B_S(\hat\Gamma)
-\max_{|S|=s}\mathcal T_{n,S}
\right|
\le C\frac{\sigma_r}{\sigma_0}\delta^3.
\end{equation}
For any fixed support $S_0$ with $|S_0|\le s$, subtracting its value from both fields
changes the right side of \eqref{eq:critical-field-maximum} by at most a
factor of two.
\end{theorem}

\begin{theorem}[Matched moment approximation]\label{thm:matched-moments}
Suppose Assumption~\ref{ass:sparse-moments} holds and $m=o(n)$.  If
\begin{equation*}
        B_* = \max_{|S|=s}B_S(\Gamma),
\end{equation*}
then
\begin{equation}\label{eq:matched-moment-approximation}
\begin{split}
        d_3\bigg(
        \frac{\max_{|S|=s}B_S(\hat\Gamma)-B_*}{a_n},
        \frac{\max_{|S|=s}\mathcal T^G_{n,S}-B_*}{a_n}
        \bigg)
        \lesssim
        \left(\frac mn\right)^{1/6}+\delta.
\end{split}
\end{equation}
For any fixed support $S_0$ with $|S_0|\le s$, put
$B_{*,S_0}=B_*-B_{S_0}(\Gamma)$.  Then
\begin{equation}\label{eq:anchored-matched-moment-approximation}
\begin{split}
 d_3\bigg(
 \frac{\max_{|S|=s}\{B_S(\hat\Gamma)-B_{S_0}(\hat\Gamma)\}
       -B_{*,S_0}}{a_n},
 \frac{\max_{|S|=s}\{\mathcal T^G_{n,S}-\mathcal T^G_{n,S_0}\}
       -B_{*,S_0}}{a_n}
 \bigg)
 \lesssim \left(\frac mn\right)^{1/6}+\delta.
\end{split}
\end{equation}
\end{theorem}

We prove the two results through sparse concentration, a uniform Taylor
expansion, and rowwise invariance of the off-diagonal quadratic field.

\begin{lemma}[Sparse hybrid moments]\label{lem:sparse-hybrid-moments}
Let $Y_i=v^{(i)}(v^{(i)})^\top-\Gamma$, and let $Y_i^G$ be independent
centered Gaussian matrices, independent of the original rows, with the covariance in
\eqref{eq:matched-moment-matrix}.  Uniformly over every deterministic
hybrid pattern $\widetilde Y_i\in\{Y_i,Y_i^G,0\}$ and each fixed
$1\le q\le6$,
\begin{equation}\label{eq:sparse-moment-concentration}
        \left\|
        \left|\frac1n\sum_{i=1}^n\widetilde Y_i\right|_s
        \right\|_{L_q}
        \lesssim\delta,
        \qquad sL=o(n).
\end{equation}
For every $t\ge0$, the corresponding tail bound is
\begin{equation}\label{eq:sparse-moment-tail}
\Prob\left(
\left|\frac1n\sum_{i=1}^n\widetilde Y_i\right|_s
>C\left\{\sqrt{\frac{m+t}{n}}+\frac{m+t}{n}\right\}
\right)
\le 2e^{-ct}.
\end{equation}
Moreover,
\begin{equation}\label{eq:one-row-moment}
        \max_i
        \left\|\lvert\widetilde Y_i\rvert_s\right\|_{L_q}
        \lesssim sL.
\end{equation}
In particular,
\begin{equation*}
        \lvert\hat\Gamma-\Gamma\rvert_s
        +\frac{\lvert G\rvert_s}{\sqrt n}
        =O_{\Prob}(\delta).
\end{equation*}
\end{lemma}

\begin{proof}
Fix a support $J$ and work on the range of $C_{JJ}$.  Every scalar
functional needed to net the three support-specific terms in
\eqref{eq:sparse-moment-size} is a centered product of two linear forms
standardized by their population variances.  It therefore has uniformly
bounded $\psi_1$ norm.  A $1/4$-net in each support-specific whitened
sphere, followed by a union over $|J|\le2s$, has logarithmic size
$O(sL)$.  Bernstein's inequality and the usual net approximation of
Euclidean and operator norms give \eqref{eq:sparse-moment-tail}, including
its linear term, for every $t\ge0$.  A matched Gaussian scalar functional
has the same variance as the corresponding original functional and is
sub-Gaussian, so the argument applies to every deterministic hybrid
pattern.

Integrating this full Bernstein tail gives
\eqref{eq:sparse-moment-concentration} for each fixed $q\le6$ when
$sL=o(n)$.  Applying the same union of whitened nets to a single original
row gives an $L_q$ bound $O(sL+q)$ for each net maximum: this follows
directly from the $\psi_1$ maximal inequality and does not require the
coordinates of $X$ to be uncorrelated.  The matched Gaussian net maximum
is $O\{\sqrt{sL+q}\}$ in $L_q$.  Since $q$ is fixed and $sL\ge1$, both
bounds imply \eqref{eq:one-row-moment}.
\end{proof}

The control-block component of \eqref{eq:sparse-moment-tail} is the
relative restricted-isometry event
\begin{equation}\label{eq:relative-rip}
        \sup_{|J|\le2s}
        \norm{C_{JJ}^{-1/2}\hat C_{JJ}C_{JJ}^{-1/2}
        -\Pi_J}_{\mathrm{op}}
        \lesssim\delta,
        \qquad \hat C=\frac{\mathbf X^\top\mathbf X}{n}.
\end{equation}
It is centered at the population covariance $C$.  The local-isotropy
condition used later instead compares $C$ with $I$; the former stabilizes
the empirical field, while the latter controls its approximation by a
separable search.

The next lemma records the deterministic part of the approximation.  Let
$DB_S(\Gamma)[H]$ and $D^2B_S(\Gamma)[H,H]$ denote the first two
derivatives of \eqref{eq:coefficient-functional} along $H$.

\begin{lemma}[Uniform moment expansion]\label{lem:uniform-expansion}
There is a constant $c>0$, depending only on the overlap constant, such
that every range-compatible $H$ with $|H|_s\le c$ satisfies
\begin{align}
\sup_{|S|\le s}
\left|B_S(\Gamma+H)-B_S(\Gamma)-DB_S(\Gamma)[H]\right|
&\lesssim \frac{\sigma_r}{\sigma_0}|H|_s^2,
\label{eq:first-order-remainder}\\
\sup_{|S|\le s}
\left|B_S(\Gamma+H)-B_S(\Gamma)-DB_S(\Gamma)[H]
-\frac12D^2B_S(\Gamma)[H,H]\right|
&\lesssim \frac{\sigma_r}{\sigma_0}|H|_s^3.
\label{eq:second-order-remainder}
\end{align}
\end{lemma}

\begin{proof}
The calculation takes place in the compressed coordinates above.  For a
fixed $S$, write
$K_S=C_{SS}^{-1}$, $\alpha=K_Sd_S$, $\gamma=K_Sh_S$, and
$\eta=\sigma_0^2-d_S^\top\alpha$.  The inverse identity
\begin{equation*}
        (C_{SS}+H_{SS})^{-1}
        =K_S-K_SH_{SS}K_S+K_SH_{SS}K_SH_{SS}K_S
        -K_SH_{SS}K_SH_{SS}K_SH_{SS}(C_{SS}+H_{SS})^{-1}
\end{equation*}
expands both Schur complements in
\eqref{eq:coefficient-functional}.  To justify these expansions uniformly
on the entire segment, put
\begin{equation*}
 E_S=C_{SS}^{-1/2}H_{SS}C_{SS}^{-1/2}.
\end{equation*}
For $0\le t\le1$, $\norm{E_S}_{\mathrm{op}}\le |H|_s<c$ and hence
\begin{equation*}
 \norm{(I+tE_S)^{-1}}_{\mathrm{op}}
 \le(1-c)^{-1}.
\end{equation*}
Moreover, $\norm{C_{SS}^{-1/2}d_S}\le\sigma_0$ and
$\norm{C_{SS}^{-1/2}H_{Sx}}\le\sigma_0|H|_s$.  Substituting these bounds
in the Schur complement gives
\begin{equation*}
 \left|(\Gamma+tH)_{xx\mathbin{\cdot}S}-\eta(S)\right|
 \le C\sigma_0^2|H|_s.
\end{equation*}
Choosing $c$ in terms of $c_{\mathrm{ov}}$ therefore keeps this denominator
above $c_{\mathrm{ov}}\sigma_0^2/2$ for every $S$ and $t$.  The unadjusted
denominator is at least $\sigma_0^2/2$ by the same choice of $c$.

The relative norm in
\eqref{eq:sparse-moment-size} bounds every linear term by its population
standard deviation, while Cauchy--Schwarz gives
$\norm{C_{SS}^{-1/2}d_S}\le\sigma_0$ and
$\norm{C_{SS}^{-1/2}h_S}\le\sigma_r$.  Overlap bounds the reciprocal
denominators by $(c_{\mathrm{ov}}\sigma_0^2)^{-1}$.  Expanding each scalar
reciprocal through second order and collecting terms of degree at least
two proves \eqref{eq:first-order-remainder}.  Retaining the quadratic
terms leaves products of degree at least three and proves
\eqref{eq:second-order-remainder}, uniformly in $S$.
\end{proof}

\begin{lemma}[Residual-field row bounds]
\label{lem:residual-field-row-bounds}
Let $V=(U^\top,X^\top)^\top$ be centered, with
$U\in\R^{d_U}$ of fixed dimension, and write
$\Gamma=\E(VV^\top)$ and $C=\E(XX^\top)$.  Let $\Pi_J$ project onto
the range of $C_{JJ}$.  Suppose
\begin{equation*}
 \norm{a^\top V}_{\psi_2}\le K\{\E(a^\top V)^2\}^{1/2}
 \qquad\text{for every }a.
\end{equation*}
Let $D_U$ be a positive diagonal matrix for which
$\norm{D_U^{-1}\E(UU^\top)D_U^{-1}}_{\mathrm{op}}\le K_U$ for a fixed
constant $K_U$, and put
\begin{equation*}
 \Xi_S=\E(X_SU^\top),\qquad
 L_S=C_{SS}^{-1}\Xi_S,\qquad
 w_S=U-L_S^\top X_S,\qquad
 R_S=\E(w_Sw_S^\top).
\end{equation*}
Call $H$ range-compatible here when
$H_{JU}=\Pi_JH_{JU}$ and $H_{JJ}=\Pi_JH_{JJ}\Pi_J$ for $|J|\le2s$.
For such a perturbation, define
\begin{equation*}
\begin{split}
 |H|_{D_U,s}={}&\norm{D_U^{-1}H_{UU}D_U^{-1}}_{\mathrm{op}}\\
 &+\sup_{|J|\le2s}\left\{
 \norm{C_{JJ}^{-1/2}H_{JU}D_U^{-1}}_{\mathrm F}
 +\norm{C_{JJ}^{-1/2}H_{JJ}C_{JJ}^{-1/2}}_{\mathrm{op}}
 \right\},\\
 \lambda_s={}&\sup_{|S|\le s}
 \norm{C_{SS}^{-1/2}\Xi_SD_U^{-1}}_{\mathrm F}.
\end{split}
\end{equation*}
Let $\varphi$ be three times continuously differentiable on a
neighborhood of the convex hull of the $R_S$.  Suppose that, for
$1\le k\le3$, every matrix $R$ in this neighborhood satisfies
\begin{equation}\label{eq:standardized-target-derivatives}
 \left|D^k\varphi(R)[D_UA_1D_U,\ldots,D_UA_kD_U]\right|
 \le\kappa_\varphi\prod_{j=1}^k\norm{A_j}_{\mathrm{op}}.
\end{equation}
For bounded scalars $c_S,c_0$, define
\begin{equation*}
 \Phi_S(A)=c_S\varphi\{\mathcal R_S^U(A)\}
 -c_0\varphi\{\mathcal R_\emptyset^U(A)\},
 \qquad
 \mathcal R_S^U(A)=A_{UU}-A_{US}A_{SS}^{-1}A_{SU}.
\end{equation*}
Then
\begin{equation}\label{eq:residual-field-derivative-bound}
 \sup_{|S|\le s}\sup_{|H|_{D_U,s}\le1}
 |D\Phi_S(\Gamma)[H]|
 \lesssim\kappa_\varphi
 \left\{\lambda_s+\sup_{|S|\le s}|c_S-c_0|\right\}.
\end{equation}
If $Y=VV^\top-\Gamma$, $Y^G$ is a centered Gaussian matrix with the
covariance of $Y$, and $H$ is independent of both $Y$ and $Y^G$, then
\begin{align}
 \sup_{|S|\le s}
 \norm{D\Phi_S(\Gamma)[Y]+D^2\Phi_S(\Gamma)[H,Y]}_{\psi_1\mid H}
 &\lesssim\kappa_\varphi\left\{
 \lambda_s+\sup_{|S|\le s}|c_S-c_0|+|H|_{D_U,s}\right\},
 \label{eq:residual-field-mixed-row-bound}\\
 \sup_{|S|\le s}\left\{
 \norm{D^2\Phi_S(\Gamma)[Y,Y]}_{L_q}
 +\norm{D^2\Phi_S(\Gamma)[Y^G,Y^G]}_{L_q}\right\}
 &\lesssim\kappa_\varphi q(s+q),\qquad q\ge2.
 \label{eq:residual-field-diagonal-bound}
\end{align}
The mixed-row bound also holds with $Y^G$ in place of $Y$, and
\begin{equation}\label{eq:residual-field-diagonal-mean}
 \E D^2\Phi_S(\Gamma)[Y,Y]
 =\E D^2\Phi_S(\Gamma)[Y^G,Y^G].
\end{equation}
If $\varphi$ ignores an entry of its matrix argument, the corresponding
terms may be omitted from $|H|_{D_U,s}$.
\end{lemma}

\begin{proof}
Fix $S$ and write
\begin{equation*}
 H_{Sw}=H_{SU}-H_{SS}L_S,
 \qquad
 H_{ww}=H_{UU}-H_{US}L_S-L_S^\top H_{SU}
 +L_S^\top H_{SS}L_S.
\end{equation*}
Differentiation on the compressed range gives
\begin{equation}\label{eq:residual-schur-derivatives}
\begin{split}
 D\mathcal R_S^U(\Gamma)[H]&=H_{ww},\\
 D^2\mathcal R_S^U(\Gamma)[H,Q]
 &=-H_{Sw}^\top C_{SS}^{-1}Q_{Sw}
 -Q_{Sw}^\top C_{SS}^{-1}H_{Sw}.
\end{split}
\end{equation}
For a centered row moment,
\begin{equation}\label{eq:row-residual-blocks}
 Y_{ww}=w_Sw_S^\top-R_S,
 \qquad Y_{Sw}=X_Sw_S^\top.
\end{equation}

Let $F_S=\nabla\varphi(R_S)$ and represent the Hessian by
$D^2\varphi(R_S)[P,Q]=\langle\mathcal H_S[P],Q\rangle$.  The incoming-row
term for one unanchored support is exactly
\begin{equation}\label{eq:residual-mixed-row-formula}
\begin{split}
 \mathcal L_{S,H}(Y)
 ={}&\langle F_S+\mathcal H_S[H_{ww}],Y_{ww}\rangle\\
 &-2\langle C_{SS}^{-1/2}H_{Sw}F_S,
 C_{SS}^{-1/2}Y_{Sw}\rangle_{\mathrm F}\\
 ={}&w_S^\top\{F_S+\mathcal H_S[H_{ww}]\}w_S
 -\operatorname{tr}[\{F_S+\mathcal H_S[H_{ww}]\}R_S]\\
 &-2X_S^\top C_{SS}^{-1}H_{Sw}F_Sw_S.
\end{split}
\end{equation}
It is a centered quadratic form in $(U,X_S)$ whose coefficient has rank
at most $3d_U$.

For the base derivative, put $K_S=(I_{d_U},-L_S^\top)^\top$.  Its coefficient
in $(U,X_S)$ is
\begin{equation*}
 c_SK_SF_SK_S^\top-c_0K_\emptyset F_\emptyset K_\emptyset^\top.
\end{equation*}
After block standardization by $\operatorname{diag}(D_U,C_{SS}^{1/2})$,
the $UU$, $UX$, and $XX$ blocks have operator norms bounded respectively
by
\begin{equation*}
 C\kappa_\varphi\{\lambda_s^2+|c_S-c_0|\},
 \qquad C\kappa_\varphi\lambda_s,
 \qquad C\kappa_\varphi\lambda_s^2.
\end{equation*}
Indeed,
$R_S-R_\emptyset=-\Xi_S^\top C_{SS}^{-1}\Xi_S$ and
$C_{SS}^{1/2}L_SD_U^{-1}=C_{SS}^{-1/2}\Xi_SD_U^{-1}$.
Positive semidefiniteness of the joint covariance and fixed $d_U$ give
$\lambda_s\le C$.  The standardized nuclear norm of the coefficient is
therefore bounded by
$C\kappa_\varphi\{\lambda_s+|c_S-c_0|\}$.  The same calculation and
\begin{equation*}
 \norm{C_{SS}^{-1/2}H_{Sw}D_U^{-1}}_{\mathrm F}
 +\norm{D_U^{-1}H_{ww}D_U^{-1}}_{\mathrm{op}}
 \le C|H|_{D_U,s}
\end{equation*}
bound the standardized nuclear norm of the mixed coefficient by
$C\kappa_\varphi|H|_{D_U,s}$.  This proves
\eqref{eq:residual-field-derivative-bound}.  Relative sub-Gaussianity and
the spectral decomposition of a fixed-rank coefficient give the required
quadratic-form bound.  To make the standardization explicit, set
$\mathsf D_S=\operatorname{diag}(D_U,C_{SS}^{1/2})$ and
$\bar\Sigma_S=\mathsf D_S^{-1}\E[(U,X_S)(U,X_S)^\top]\mathsf D_S^{-1}$.
Positive semidefiniteness and the assumed target-block bound imply
$\norm{\bar\Sigma_S}_{\mathrm{op}}\le C$.  Hence, for every symmetric
coefficient $A$,
\begin{equation*}
 \norm{\bar\Sigma_S^{1/2}(\mathsf D_SA\mathsf D_S)
       \bar\Sigma_S^{1/2}}_*
 \le C\norm{\mathsf D_SA\mathsf D_S}_*.
\end{equation*}
Applying the relative sub-Gaussian bound after this standardization, with
$z=(U^\top,X_S^\top)^\top$ and $\Sigma_z=\E(zz^\top)$, yields
\begin{equation*}
 \norm{z^\top Az-\E(z^\top Az)}_{\psi_1}
 \le CK^2\norm{\Sigma_z^{1/2}A\Sigma_z^{1/2}}_*,
\end{equation*}
which proves \eqref{eq:residual-field-mixed-row-bound}.  The corresponding
functional of $Y^G$ is Gaussian with the same variance and obeys the same
bound.

For the diagonal term, the chain rule gives the exact identity
\begin{equation}\label{eq:residual-diagonal-formula}
\begin{split}
 D^2(\varphi\circ\mathcal R_S^U)(\Gamma)[Y,Y]
 ={}&D^2\varphi(R_S)[Y_{ww},Y_{ww}]\\
 &-2D\varphi(R_S)
 [Y_{Sw}^\top C_{SS}^{-1}Y_{Sw}].
\end{split}
\end{equation}
For an original row,
\begin{equation*}
 Y_{Sw}^\top C_{SS}^{-1}Y_{Sw}
 =(X_S^\top C_{SS}^{-1}X_S)w_Sw_S^\top.
\end{equation*}
Relative sub-Gaussianity gives
\begin{equation*}
 \norm{X_S^\top C_{SS}^{-1}X_S}_{L_{2q}}\lesssim s+q,
 \qquad
 \norm{D_U^{-1}w_Sw_S^\top D_U^{-1}}_{L_{2q}}\lesssim q.
\end{equation*}
Together with \eqref{eq:standardized-target-derivatives}, this proves the
first term in \eqref{eq:residual-field-diagonal-bound}.

For the matched Gaussian row, put
\begin{equation*}
 \mathcal Z_S=C_{SS}^{-1/2}(Y^G)_{Sw}D_U^{-1}.
\end{equation*}
If $\norm A_{\mathrm F}=1$, covariance matching and a singular-value
decomposition of the $|S|\times d_U$ matrix $A$ give
\begin{equation*}
 \Var\langle A,\mathcal Z_S\rangle
 =\Var\langle A,C_{SS}^{-1/2}X_Sw_S^\top D_U^{-1}\rangle
 \le C.
\end{equation*}
Thus the covariance operator of $\operatorname{vec}(\mathcal Z_S)$ is
bounded and its trace is $O(sd_U)$.  Gaussian quadratic-form moments yield
$\norm{\mathcal Z_S^\top\mathcal Z_S}_{L_q}\lesssim s+q$.
The fixed-dimensional block
$D_U^{-1}(Y^G)_{ww}D_U^{-1}$ has $L_{2q}$ norm $O(\sqrt q)$, so
\eqref{eq:residual-diagonal-formula} proves the Gaussian part of
\eqref{eq:residual-field-diagonal-bound}.  Finally,
\eqref{eq:residual-field-diagonal-mean} follows because the bilinear form
$D^2\Phi_S(\Gamma)$ is deterministic and $Y,Y^G$ have the same
covariance.
\end{proof}

\begin{proof}[Proof of Theorem~\ref{thm:critical-empirical-expansion}]
Lemma~\ref{lem:sparse-hybrid-moments}, with $t=m$, gives
$|\hat\Gamma-\Gamma|_s\le C\delta$ outside an event of probability
$Ce^{-cm}$.  Support compression makes the same statement valid when a
population block is singular.  On this event the compressed empirical
Gram blocks are invertible, and overlap gives
$\hat\Gamma_{xx\mathbin{\cdot}S}\ge
c_{\mathrm{ov}}\sigma_0^2/2$ uniformly over $|S|\le s$.  Thus every
searched coefficient is defined.  Lemma~\ref{lem:uniform-expansion}
gives the two displays.  The maximum statement follows from
$|\max_S a_S-\max_S b_S|\le\max_S|a_S-b_S|$.
\end{proof}

\begin{proof}[Proof of Theorem~\ref{thm:matched-moments}]
Apply Lemma~\ref{lem:residual-field-row-bounds} with
\begin{equation*}
 U=(x,r)^\top,\qquad
 D_U=\operatorname{diag}(\sigma_0,\sigma_r),\qquad
 \varphi(R)=R_{12}/R_{11},\qquad c_S=c_0=1.
\end{equation*}
Overlap gives \eqref{eq:standardized-target-derivatives} with
$\kappa_\varphi=C\sigma_r/\sigma_0$.  Moreover,
\begin{equation*}
 \lambda_s
 \le\sqrt{\phi_s}+\frac{\widetilde\Psi_s}{\sigma_r}.
\end{equation*}
The function $\varphi$ ignores the $rr$ entry, so the reduced target norm
in the final sentence of the lemma is bounded by $C|H|_s$.  Consequently,
for $H$ independent of a centered row moment $Y$ and $q\ge2$,
\begin{align}
 \sup_{|S|=s}
 \norm{DB_S(\Gamma)[Y]+D^2B_S(\Gamma)[H,Y]}_{\psi_1\mid H}
 &\lesssim\frac1{\sigma_0}
 \{\sigma_r\sqrt{\phi_s}+\widetilde\Psi_s
 +\sigma_r|H|_s\},
 \label{eq:intrinsic-linear-row-bound}\\
 \sup_{|S|=s}
 \norm{D^2B_S(\Gamma)[Y,Y]}_{L_q}
 &\lesssim\frac{\sigma_r}{\sigma_0}q(s+q).
 \label{eq:intrinsic-diagonal-row-bound}
\end{align}
Both bounds hold for $Y^G$, and
\begin{equation}\label{eq:diagonal-mean-matching}
        \E D^2B_S(\Gamma)[Y,Y]
        =\E D^2B_S(\Gamma)[Y^G,Y^G],
\end{equation}
by \eqref{eq:residual-field-diagonal-mean}.

We now compare the two Taylor fields.  For either row law, put
\begin{equation*}
        \mathsf{Diag}_{n,S}
        =\frac1{2n^2}\sum_{i=1}^n
        \{D^2B_S(\Gamma)[Y_i,Y_i]
        -\E D^2B_S(\Gamma)[Y,Y]\}.
\end{equation*}
The quadratic Taylor term decomposes exactly as
\begin{equation}\label{eq:taylor-diagonal-split}
\begin{split}
 \frac12D^2B_S(\Gamma)\left[\frac1n\sum_iY_i,
                         \frac1n\sum_iY_i\right]
 ={}&\frac1{n^2}\sum_{i<j}D^2B_S(\Gamma)[Y_i,Y_j]
 +\frac1{2n}\E D^2B_S(\Gamma)[Y,Y]+\mathsf{Diag}_{n,S}.
\end{split}
\end{equation}
For centered independent $Z_i$ with the moment growth in
\eqref{eq:intrinsic-diagonal-row-bound}, the independent-sum inequality
\begin{equation*}
 \norm{\textstyle\sum_{i=1}^n Z_i}_{L_\nu}
 \lesssim
 \sup_{2\le q\le\nu}
 \frac\nu q\left(\frac n\nu\right)^{1/q}\norm{Z_1}_{L_q}
 \lesssim \frac{\sigma_r}{\sigma_0}
 \{s\sqrt{n\nu}+\nu(s+\nu)\}
\end{equation*}
follows by symmetrization and the standard moment bound for sums.  A
union over at most $e^m$ supports then gives, for $\nu=m+u\lesssim n$,
\begin{equation}\label{eq:diagonal-field-concentration}
 \Prob\left[
 \max_{|S|=s}|\mathsf{Diag}_{n,S}|
 >C\frac{\sigma_r}{\sigma_0}
 \left\{\frac{s\sqrt\nu}{n^{3/2}}
 +\frac{\nu(s+\nu)}{n^2}\right\}
 \right]\le e^{-c\nu}.
\end{equation}
Since there are at most $e^m$ supports, the moment bound preceding
\eqref{eq:diagonal-field-concentration}, applied with
$\nu_0=m\vee2$, also gives directly
\begin{equation*}
 \left\|\max_{|S|=s}|\mathsf{Diag}_{n,S}|\right\|_{L_{\nu_0}}
 \lesssim \frac{\sigma_r}{\sigma_0}
 \left\{\frac{s\sqrt{\nu_0}}{n^{3/2}}
 +\frac{\nu_0(s+\nu_0)}{n^2}\right\}.
\end{equation*}
Here the factor from the maximum is at most
$e^{m/\nu_0}\le e$.  Because
$a_n\ge(\sigma_r/\sigma_0)m/n$, $s\le m$, and
$\nu_0\asymp m$, division by $a_n$ bounds the right side by
\begin{equation*}
 C\left\{\frac{s}{\sqrt{nm}}+\frac{s+m}{n}\right\}
 \le C\delta.
\end{equation*}
Thus the centered diagonal field is $O(\delta)$ after normalization in
$L_1$, and hence also for the truncated loss.  Its retained mean agrees
for the two row laws by \eqref{eq:diagonal-mean-matching}.

It remains to replace the off-diagonal field.  At the $i$th step of a
row-by-row replacement, let
\begin{equation*}
        H_i=\frac1n\sum_{j\ne i}\widetilde Y_j.
\end{equation*}
Conditional on $H_i$, inserting a row $Y$ changes the normalized
support value by the linear functional
\begin{equation*}
        \Delta_{i,S}(Y)
        =\frac{DB_S(\Gamma)[Y]+D^2B_S(\Gamma)[H_i,Y]}
        {na_n}.
\end{equation*}
The original and Gaussian increments have the same conditional mean and
covariance.  Equation~\eqref{eq:intrinsic-linear-row-bound} and the
definition of $a_n$ give
\begin{equation}\label{eq:off-diagonal-row-bound}
        \sup_{|S|=s}\norm{\Delta_{i,S}(Y)}_{\psi_1\mid H_i}
        +\sup_{|S|=s}\norm{\Delta_{i,S}(Y^G)}_{\psi_1\mid H_i}
        \lesssim\frac{\Lambda_i}{\sqrt{mn}},
        \qquad
        \Lambda_i=1+\frac{|H_i|_s}{\delta}.
\end{equation}
The bound also holds for signed combinations with bounded total
variation.  Lemma~\ref{lem:sparse-hybrid-moments} implies bounded fixed
moments of $\Lambda_i$.  Applying its exponential tail with
$u\asymp n^{2/3}m^{1/3}=o(n)$ also gives
\begin{equation*}
        \sum_{i=1}^n\Prob\{\Lambda_i>c(n/m)^{1/3}\}
        =o\{(m/n)^{1/6}\}.
\end{equation*}

Smooth the maximum with inverse temperature $\lambda m$.  Its endpoint
error is at most $1/\lambda$.  A second-order Taylor expansion in
$\Delta_i$ cancels conditionally through order two.  Let $\pi_S$ be the
Gibbs weights before the insertion and $\pi_S(\theta)$ those after
inserting $\theta\Delta_i$.  Convexity of the log-partition function gives,
for $k\le3$,
\begin{equation*}
 \prod_{a=1}^k\pi_{S_a}(\theta)
 \le \prod_{a=1}^k\pi_{S_a}
 \exp\!\left[\theta\lambda m\left\{
 \sum_{a=1}^k\Delta_{i,S_a}
 -k\sum_T\pi_T\Delta_{i,T}\right\}\right].
\end{equation*}
The expression in braces is a signed combination with conditional
$\psi_1$ norm at most $C\Lambda_i/\sqrt{mn}$.  On
$\Lambda_i\le c(n/m)^{1/3}$, the choice
$\lambda=(n/m)^{1/6}$ makes its tilted exponential moment bounded, since
$\lambda\sqrt{m/n}\,\Lambda_i\le c$.  Conditional H\"older inequalities
and \eqref{eq:off-diagonal-row-bound} then bound the one-row third-order
remainder by
\begin{equation*}
 C\{1+\lambda m+(\lambda m)^2\}
 \frac{\Lambda_i^3}{(mn)^{3/2}}.
\end{equation*}
After summing the rows and adding the smoothing error, the bound is
\begin{equation*}
        C\left\{
        \frac1\lambda
        +
        \frac1{m^{3/2}\sqrt n}
        +\frac{\lambda}{\sqrt{mn}}
        +\lambda^2\sqrt{\frac mn}\right\}.
\end{equation*}
Taking $\lambda=(n/m)^{1/6}$ proves an
$O\{(m/n)^{1/6}\}$ comparison between the off-diagonal fields.

Finally, take $t=c_0n$ in \eqref{eq:sparse-moment-tail}, with $c_0>0$
small enough that both endpoints lie in the Taylor neighborhood outside
an event of probability $Ce^{-cn}$.  On this event,
Lemma~\ref{lem:uniform-expansion} and the fixed-moment form of sparse
concentration bound the cubic remainder at both endpoints by
$C(\sigma_r/\sigma_0)\delta^3$ in truncated first moment.  Explicitly,
the full tail gives
\begin{equation*}
 \E\left[1\wedge
 \frac{(\sigma_r/\sigma_0)|H|_s^3}{a_n}\right]
 \le C\delta,
\end{equation*}
because $a_n\ge(\sigma_r/\sigma_0)\delta^2$; bounded test functions
absorb the exceptional event.
Together with \eqref{eq:diagonal-field-concentration}, this proves
\eqref{eq:matched-moment-approximation}.  Subtracting one fixed support
changes each bound by at most a constant factor.
\end{proof}

For the linear specialization
\begin{equation}\label{eq:model}
        y^{(i)}=\beta_0x^{(i)}+(X^{(i)})^\top\beta+\varepsilon^{(i)},
        \qquad \E[(\varepsilon^{(i)})^2]=\sigma_\varepsilon^2,
\end{equation}
let the centered noise $\varepsilon^{(i)}$ be independent of
$(x^{(i)},X^{(i)})$.  Then
\begin{equation}\label{eq:score-translation}
\begin{split}
        h&=C\beta-\frac{d^\top\beta}{\sigma_0^2}d,
        \qquad
        \sigma_r^2
        =\beta^\top C\beta+\sigma_\varepsilon^2
        -\frac{(d^\top\beta)^2}{\sigma_0^2},\\
        \widetilde\Psi_k
        &\le \Psi_k
        +\frac{|d^\top\beta|}{\sigma_0^2}\norm{d}_{C,k},
        \qquad
        \Psi_k=\norm{C\beta}_{C,k}.
\end{split}
\end{equation}
Thus $\widetilde\Psi_k$ is the residual-score size, whereas $\Psi_k$ is
the outcome-score size.  Under random assignment, $d=0$,
$\beta_\emptyset^*=\beta_0$, and the two coincide.
Every statement below involving $\beta_0$, $\beta$, $\Psi_s$, or $\varepsilon$ uses
the linear specialization \eqref{eq:model}.

The first derivative has the regression influence form.  In a population
row, put
\begin{equation}\label{eq:influence-residuals}
\begin{aligned}
        \alpha_S&=C_{SS}^{-1}d_S,
        &u_S&=x-X_S^\top\alpha_S,\\
        \gamma_S&=C_{SS}^{-1}h_S,
        &e_S&=r-X_S^\top\gamma_S-B_S(\Gamma)u_S.
\end{aligned}
\end{equation}
Then
\begin{equation}\label{eq:influence-function}
        DB_S(\Gamma)[vv^\top-\Gamma]
        =\frac{u_Se_S}{\eta(S)}-\frac{xr}{\sigma_0^2}.
\end{equation}
Thus the Taylor field retains the full observational influence process;
sparse treatment alignment reduces it further to the score field below.

\subsection{Score-product reduction}

The general field simplifies when the treatment has little population
projection on the controls.  For \(a,b\in\R^p\), define
\begin{equation}\label{eq:chaos-functional}
        Q_S(a,b)=-a_S^\top C_{SS}^{-1}b_S,
        \qquad
        T_C(a,b)=\max_{|S|=s}Q_S(a,b).
\end{equation}
The score pairs in the matched and empirical moment fields are
\begin{equation}\label{eq:matched-scores}
        Z_1=\sqrt n\,d+G_{Xx},
        \qquad
        Z_2=\sqrt n\,h+G_{Xr},
\end{equation}
and
\begin{equation}\label{eq:empirical-scores}
        Z_{1n}=\frac{\mathbf X^\top\mathbf x}{\sqrt n},
        \qquad
        Z_{2n}=\frac{\mathbf X^\top\mathbf r}{\sqrt n}.
\end{equation}
The centered matched pair has the covariance of the row scores
\(X^{(i)}x^{(i)}-d\) and \(X^{(i)}r^{(i)}-h\).  Put
\begin{equation}\label{eq:local-treatment-scale}
        a_n^{\mathrm{score}}
        =\frac{\widetilde\Psi_s}{\sigma_0}\delta
        +\frac{\sigma_r}{\sigma_0}\delta^2.
\end{equation}

\begin{theorem}[Sparse-treatment score approximation]
\label{thm:chaos-degeneration}
Suppose Assumption~\ref{ass:sparse-moments} holds, $\delta\le\delta_0$, and
\begin{equation}\label{eq:local-treatment}
        \sqrt{\phi_{2s}}\lesssim\delta.
\end{equation}
No restriction is imposed on \(h\).  Then
\begin{align}
\sup_{|S|\le s}\left|
        \mathcal T^G_{n,S}
        -\frac{Q_S(Z_1,Z_2)}{n\sigma_0^2}\right|
&=O_{\Prob}(\delta a_n^{\mathrm{score}}),
\label{eq:gaussian-chaos-degeneration}\\
\sup_{|S|\le s}\left|
        B_S(\hat\Gamma)
        -\frac{Q_S(Z_{1n},Z_{2n})}{n\sigma_0^2}\right|
&=O_{\Prob}(\delta a_n^{\mathrm{score}}).
\label{eq:empirical-chaos-degeneration}
\end{align}
Both remainders obey this scale outside an event of probability
\(Ce^{-cm}\).  If, moreover, \(m=o(n)\), then
\begin{equation}\label{eq:local-treatment-stochastic-approximation}
        d_3\bigg(
        \frac{\max_{|S|=s}
        (\hat\beta_{0,S}-\hat\beta_{0,\emptyset})}
        {a_n^{\mathrm{score}}},
        \frac{T_C(Z_1,Z_2)}
        {n\sigma_0^2a_n^{\mathrm{score}}}
        \bigg)
        \lesssim
        \left(\frac{m}{n}\right)^{1/6}+\delta.
\end{equation}
\end{theorem}

\begin{proof}
For a range-compatible perturbation \(H\), write
\[
        a_S=d_S+H_{Sx},\qquad b_S=h_S+H_{Sr},\qquad
        \widetilde\sigma_0^2=\sigma_0^2+H_{xx},\qquad
        \widetilde C_S=C_{SS}+H_{SS}.
\]
Direct substitution into the coefficient functional gives
\begin{equation}\label{eq:exact-local-treatment-perturbation}
        B_S(\Gamma+H)
        =\frac{
        -a_S^\top\widetilde C_S^{-1}b_S
        +(H_{xr}/\widetilde\sigma_0^2)
        a_S^\top\widetilde C_S^{-1}a_S}
        {\widetilde\sigma_0^2
        -a_S^\top\widetilde C_S^{-1}a_S}.
\end{equation}
The leading term is
\(-a_S^\top C_{SS}^{-1}b_S/\sigma_0^2\).  Relative inverse bounds make
the uniform difference, on \(|H|_s\le c\), at most a constant times
\begin{equation}\label{eq:score-reduction-deterministic}
\begin{split}
 \frac{\sigma_r}{\sigma_0}
 &\left(\sqrt{\phi_{2s}}+|H|_s\right)
 \left\{\frac{\widetilde\Psi_s}{\sigma_r}+|H|_s\right\}
 \left\{|H|_s+\left(\sqrt{\phi_{2s}}+|H|_s\right)^2\right\}\\
 &+\frac{\sigma_r}{\sigma_0}
 \left(\sqrt{\phi_{2s}}+|H|_s\right)^2|H|_s .
\end{split}
\end{equation}
Take \(H=G/\sqrt n\) and \(H=\hat\Gamma-\Gamma\).  The sparse block
bounds from Lemma~\ref{lem:sparse-hybrid-moments}, together with
\(\sqrt{\phi_{2s}}\lesssim\delta\), turn
\eqref{eq:score-reduction-deterministic} into
\(\delta a_n^{\mathrm{score}}\).  Lemma~\ref{lem:uniform-expansion}
replaces $B_S(\Gamma+G/\sqrt n)$ by $\mathcal T^G_{n,S}$ at the smaller
cubic order.  More precisely, put
\begin{equation*}
\begin{split}
 R_n^G&=\sup_{|S|\le s}\left|
 \mathcal T^G_{n,S}-\frac{Q_S(Z_1,Z_2)}{n\sigma_0^2}\right|,\\
 R_n^E&=\sup_{|S|\le s}\left|
 B_S(\hat\Gamma)-\frac{Q_S(Z_{1n},Z_{2n})}{n\sigma_0^2}\right|.
\end{split}
\end{equation*}
The fixed-moment and full-tail parts of
Lemma~\ref{lem:sparse-hybrid-moments}, applied to the polynomial in
\eqref{eq:score-reduction-deterministic}, give
\begin{equation}\label{eq:score-reduction-truncated-error}
 \E\left(1\wedge\frac{R_n^G}{a_n^{\mathrm{score}}}\right)
 +\E\left(1\wedge\frac{R_n^E}{a_n^{\mathrm{score}}}\right)
 \lesssim\delta.
\end{equation}
On $|H|_s\le C\delta$, the same calculation gives the asserted
high-probability bounds.  This proves the first two claims.

Under \eqref{eq:local-treatment}, the scales in
\eqref{eq:mean-sizes} and \eqref{eq:local-treatment-scale} satisfy
$a_n^{\mathrm{score}}\le a_n\le C a_n^{\mathrm{score}}$.  Theorem
\ref{thm:matched-moments} therefore compares the empirical and matched
Taylor maxima at the score scale.  The first claim replaces the latter by
$T_C(Z_1,Z_2)/(n\sigma_0^2)$; equation
\eqref{eq:score-reduction-truncated-error} bounds the resulting change in
$d_3$ by $C\delta$.  The triangle inequality proves
\eqref{eq:local-treatment-stochastic-approximation}.
\end{proof}

The two terms in \(a_n^{\mathrm{score}}\) are respectively the linear
outcome-score and bilinear-chaos scales.  The theorem also gives, with
probability at least \(1-Ce^{-cm}\),
\begin{equation}\label{eq:local-treatment-upper-scale}
        \max_{|S|=s}
        (\hat\beta_{0,S}-\hat\beta_{0,\emptyset})
        \le Ca_n^{\mathrm{score}},
        \qquad \delta\le\delta_0,
\end{equation}
Indeed, on the sparse score event,
\begin{equation*}
 \norm{Z_{1n}}_{C,s}\le C\sigma_0\sqrt m,
 \qquad
 \norm{Z_{2n}}_{C,s}
 \le C\{\sqrt n\,\widetilde\Psi_s+\sigma_r\sqrt m\}.
\end{equation*}
Together with \(T_C(a,b)\le\norm{a}_{C,s}\norm{b}_{C,s}\) and
\eqref{eq:empirical-chaos-degeneration}, these bounds prove
\eqref{eq:local-treatment-upper-scale}.  Under random
assignment this is equivalent to
\[
        \frac{\delta}{\sigma_0}
        (\Psi_s^2+\sigma_r^2\delta^2)^{1/2}.
\]
A matching scalar lower bound cannot depend on \(\Psi_s\) alone: it also
requires enough separated supports carrying that energy.
\section{Certified product search}

\subsection{Conditional products and projection certificate}

A vector $g$ is score-compatible when
$g_J\in\operatorname{range}(C_{JJ})$ for every
$|J|\le2s$; all population, empirical, and matched score vectors have
this property.  If $j\notin S$ and
$C_{jj}-C_{jS}C_{SS}^{-1}C_{Sj}>0$, define its standardized residual score
\begin{equation}\label{eq:conditional-score}
        g_{j\mid S}
        =\frac{g_j-C_{jS}C_{SS}^{-1}g_S}
        {\{C_{jj}-C_{jS}C_{SS}^{-1}C_{Sj}\}^{1/2}},
\end{equation}
and put $g_{j\mid S}=0$ when the residual variance is zero.
Compatibility makes the residual numerator vanish in the latter case.
Block inversion gives
\begin{equation}\label{eq:conditional-product}
        Q_{S\cup\{j\}}(a,b)-Q_S(a,b)
        =-a_{j\mid S}b_{j\mid S}.
\end{equation}
When the residual variance is zero, both sides of
\eqref{eq:conditional-product} vanish.
The literal extension of top-$s$ sorting starts at $J_0=\emptyset$ and
selects the largest current product $-a_{j\mid J_{t-1}}b_{j\mid
J_{t-1}}$ at each step.  It is exact when $C$ is diagonal, but under
correlation an early choice changes every later product.  The
certifiable construction below instead selects a polarized projection
gain and bounds the regret of its resulting product path.

For disjoint $S,R\subseteq[p]$, put
\begin{equation}\label{eq:conditional-control-block}
        C_{R\mid S}
        =C_{RR}-C_{RS}C_{SS}^{-1}C_{SR}.
\end{equation}
The scale-invariant projection ratio is
\begin{equation}\label{eq:residual-ratio}
        \mu_s(C)
        =\inf_{S,R,u}
        \frac{u^\top\operatorname{diag}(C_{R\mid S})^{-1}u}
        {u^\top C_{R\mid S}^{-1}u},
\end{equation}
where the infimum is over disjoint sets with $|S|\le s$,
$1\le|R|\le s$, and nonzero
$u\in\operatorname{range}(C_{R\mid S})$; pairs with
$C_{R\mid S}=0$ are omitted, and an empty infimum is one.  Equivalently,
\eqref{eq:residual-ratio} is the infimum of the least positive
eigenvalues of the conditional correlation matrices obtained from
\eqref{eq:conditional-control-block} after zero-variance coordinates are
removed.  This definition allows singular control blocks.

This parameter is the largest covariance-only constant in the
joint-to-coordinate projection inequality used below.  Indeed, for every
admissible $u$, the compatible score
$g=C_{\cdot R}C_{R\mid S}^{-1}u$ has conditional residual $u$.  Thus every
direction in the infimum is attained.  The
parameter also connects directly to sparse restricted conditioning.
When
every $C_{JJ}$ with $|J|\le2s$ is nonsingular, put
\begin{equation*}
        \kappa_s(C)=\sup_{|J|\le2s}
        \frac{\lambda_{\max}(C_{JJ})}{\lambda_{\min}(C_{JJ})}.
\end{equation*}
Because the controls have unit variance, this is the sparse restricted
condition number.  For any admissible $S,R$,
$C_{R\mid S}^{-1}$ is a principal block of
$C_{S\cup R,S\cup R}^{-1}$.  Hence the eigenvalues of $C_{R\mid S}$ lie
between those of its parent block.  Its diagonal entries are bounded by
the same upper eigenvalue, so diagonal normalization gives
\begin{equation}\label{eq:projection-ratio-conditioning}
        \mu_s(C)\ge\kappa_s(C)^{-1}.
\end{equation}

Fix a supplied lower bound $0<\mu\le\mu_s(C)$.  For positive
$\sigma_a,\sigma_b$, define the active and passive contrasts
\begin{equation}\label{eq:score-polarization}
        z_{\mathrm a}=\frac1{\sqrt2}
        \left(\frac a{\sigma_a}-\frac b{\sigma_b}\right),
        \qquad
        z_{\mathrm p}=\frac1{\sqrt2}
        \left(\frac a{\sigma_a}+\frac b{\sigma_b}\right).
\end{equation}
Starting from $J_0=\emptyset$, forward projection chooses
\begin{equation}\label{eq:forward-projection-path}
        j_t\in\mathop{\arg\max}_{j\notin J_{t-1}}
        (z_{\mathrm a})_{j\mid J_{t-1}}^2,
        \qquad
        J_t=J_{t-1}\cup\{j_t\},
        \qquad t=1,\ldots,s.
\end{equation}
For a finite family, braces with subscript $(\ell)$ denote its
$\ell$th largest member.  Along this path, define
\begin{equation}\label{eq:projection-certificate}
\begin{split}
        \mathcal U_{C,\mu}(g;J_\bullet)
        =\min_{0\le t\le s}\bigg\{
        g_{J_t}^\top C_{J_tJ_t}^{-1}g_{J_t}
        +\frac1\mu
        \sum_{\ell=1}^{s\wedge(p-t)}
        \{g_{j\mid J_t}^2:j\notin J_t\}_{(\ell)}
        \bigg\}.
\end{split}
\end{equation}
Write the corresponding product-gap certificate as
\begin{equation}\label{eq:product-gap-certificate}
\begin{split}
        \mathcal D_{C,\mu}(z_{\mathrm a},z_{\mathrm p};J_\bullet)
        ={}&\mathcal U_{C,\mu}(z_{\mathrm a};J_\bullet)
        -(z_{\mathrm a})_{J_s}^\top C_{J_sJ_s}^{-1}(z_{\mathrm a})_{J_s}\\
        &+(z_{\mathrm p})_{J_s}^\top C_{J_sJ_s}^{-1}(z_{\mathrm p})_{J_s}.
\end{split}
\end{equation}

\begin{theorem}[Forward product-search guarantee]
\label{thm:projection-products}
Let $a,b$ be score-compatible, form $z_{\mathrm a},z_{\mathrm p}$ by
\eqref{eq:score-polarization}, and let $J_\bullet$ follow
\eqref{eq:forward-projection-path}.  For every score-compatible $g$,
\begin{equation}\label{eq:projection-upper-bound}
        \mathcal U_{C,\mu}(g;J_\bullet)\ge\norm{g}_{C,s}^2.
\end{equation}
The active gap obeys
\begin{equation}\label{eq:forward-projection-gap}
\begin{split}
        0&\le\norm{z_{\mathrm a}}_{C,s}^2
        -(z_{\mathrm a})_{J_t}^\top C_{J_tJ_t}^{-1}(z_{\mathrm a})_{J_t}\\
        &\le\left(1-\frac\mu s\right)^t\norm{z_{\mathrm a}}_{C,s}^2,
        \qquad 0\le t\le s.
\end{split}
\end{equation}
The returned conditional-product sum satisfies
\begin{equation}\label{eq:forward-product-certificate}
\begin{split}
        0&\le T_C(a,b)-Q_{J_s}(a,b)\\
        &\le\frac{\sigma_a\sigma_b}{2}
        \mathcal D_{C,\mu}(z_{\mathrm a},z_{\mathrm p};J_\bullet).
\end{split}
\end{equation}
Incremental factorization computes the path and the certificates in
$O(ps^2)$ arithmetic operations and $O(ps)$ working memory beyond access
to $C$.
\end{theorem}

\begin{proof}
For disjoint $S,R$, write
\begin{equation*}
        u=g_R-C_{RS}C_{SS}^{-1}g_S.
\end{equation*}
Block inversion on the compressed range and
\eqref{eq:residual-ratio} give
\begin{equation}\label{eq:joint-projection-gain}
\begin{split}
        g_{S\cup R}^\top C_{S\cup R,S\cup R}^{-1}g_{S\cup R}
        -g_S^\top C_{SS}^{-1}g_S
        &=u^\top C_{R\mid S}^{-1}u,\\
        \sum_{j\in R}g_{j\mid S}^2
        &=u^\top\operatorname{diag}(C_{R\mid S})^{-1}u\\
        &\ge\mu u^\top C_{R\mid S}^{-1}u.
\end{split}
\end{equation}
Let $A$ maximize $g_A^\top C_{AA}^{-1}g_A$ over $|A|=s$ and take
$R=A\setminus J_t$.  Projection energy is monotone, so
\eqref{eq:joint-projection-gain} yields
\begin{equation*}
\begin{split}
        \norm{g}_{C,s}^2
        -g_{J_t}^\top C_{J_tJ_t}^{-1}g_{J_t}
        \le\frac1\mu
        \sum_{\ell=1}^{s\wedge(p-t)}
        \{g_{j\mid J_t}^2:j\notin J_t\}_{(\ell)}.
\end{split}
\end{equation*}
This proves \eqref{eq:projection-upper-bound}.  With $g=z_{\mathrm a}$, the largest
residual gain is at least $\mu/s$ times the current active gap, and
iteration proves \eqref{eq:forward-projection-gap}.

Residualization commutes with \eqref{eq:score-polarization}, and
polarization gives
\begin{equation}\label{eq:conditional-product-polarization}
\begin{split}
        -a_{j\mid S}b_{j\mid S}
        &=\frac{\sigma_a\sigma_b}{2}
        \{(z_{\mathrm a})_{j\mid S}^2-(z_{\mathrm p})_{j\mid S}^2\},\\
        Q_S(a,b)
        &=\frac{\sigma_a\sigma_b}{2}\left\{
        (z_{\mathrm a})_{S}^\top C_{SS}^{-1}(z_{\mathrm a})_{S}
        -(z_{\mathrm p})_{S}^\top C_{SS}^{-1}(z_{\mathrm p})_{S}\right\}.
\end{split}
\end{equation}
Consequently,
\begin{equation*}
\begin{split}
        &\{T_C(a,b)-Q_{J_s}(a,b)\}
        +\left\{\frac{\sigma_a\sigma_b}{2}
        \norm{z_{\mathrm a}}_{C,s}^2-T_C(a,b)\right\}\\
        &\qquad=\frac{\sigma_a\sigma_b}{2}\left\{
        \norm{z_{\mathrm a}}_{C,s}^2
        -(z_{\mathrm a})_{J_s}^\top C_{J_sJ_s}^{-1}(z_{\mathrm a})_{J_s}
        +(z_{\mathrm p})_{J_s}^\top C_{J_sJ_s}^{-1}(z_{\mathrm p})_{J_s}\right\}.
\end{split}
\end{equation*}
Both terms on the left are nonnegative.  Dropping the second and applying
\eqref{eq:projection-upper-bound} proves
\eqref{eq:forward-product-certificate}.

For the computational claim, maintain a rank-revealing Cholesky or QR
factorization of $C_{J_tJ_t}$ and the residual score numerators for all
remaining coordinates.  Updating one candidate at stage $t$ costs
$O(t)$ operations, and linear-time selection finds the largest gain and
the $s\wedge(p-t)$ largest certificate terms.  A zero pivot has zero
compatible residual score, so its numerical factor update is skipped;
the fixed tie rule may still append that coordinate to the path.  Summing $O(pt)$ over
$t\le s$ gives $O(ps^2)$ operations; the candidate-by-selected-coordinate
updates require $O(ps)$ working memory.
\end{proof}

The path does not require $\mu$; the supplied lower bound makes its
realized regret certificate valid.  The best sparse value of $\mu$ can
be combinatorial, but a structural or conservative spectral bound is
enough.  Relative restricted isometry transfers such a population bound
to the empirical Gram matrix.

\begin{lemma}[Relative RIP transfer]
\label{lem:projection-rip-transfer}
Suppose positive semidefinite $C$ and $\widetilde C$ satisfy
\begin{equation}\label{eq:relative-rip-loewner}
        (1-\epsilon)C_{JJ}
        \preceq\widetilde C_{JJ}
        \preceq(1+\epsilon)C_{JJ},
        \qquad |J|\le2s,
\end{equation}
for some $0\le\epsilon<1$.  Then
\begin{equation}\label{eq:projection-rip-transfer}
        \mu_s(\widetilde C)
        \ge\frac{1-\epsilon}{1+\epsilon}\mu_s(C).
\end{equation}
\end{lemma}

\begin{proof}
The Schur complement has the variational representation
\begin{equation*}
        u^\top C_{R\mid S}u
        =\inf_w
        \begin{pmatrix}w\\u\end{pmatrix}^{\!\top}
        C_{S\cup R,S\cup R}
        \begin{pmatrix}w\\u\end{pmatrix}.
\end{equation*}
Thus \eqref{eq:relative-rip-loewner} gives
\begin{equation*}
        (1-\epsilon)C_{R\mid S}
        \preceq\widetilde C_{R\mid S}
        \preceq(1+\epsilon)C_{R\mid S}.
\end{equation*}
The two conditional matrices have the same range and zero diagonal
coordinates.  On their common range, their inverse quadratic forms obey
the reversed inequalities, while the diagonal inverse forms obey the
corresponding coordinatewise inequalities.  Hence every ratio in
\eqref{eq:residual-ratio} for $\widetilde C$ is at least
$(1-\epsilon)/(1+\epsilon)$ times its counterpart for $C$.  Taking the
infimum proves \eqref{eq:projection-rip-transfer}.
\end{proof}

In particular, \eqref{eq:relative-rip} and
\eqref{eq:projection-ratio-conditioning} give
\begin{equation}\label{eq:empirical-projection-ratio}
        \mu_s(\hat C)
        \ge\{1-O_{\Prob}(\delta)\}\mu_s(C)
        \ge\frac{1-O_{\Prob}(\delta)}{\kappa_s(C)}
\end{equation}
when the sparse population blocks are nonsingular.  Thus
$\kappa_s(C)$ is a familiar sufficient condition, whereas $\mu_s(C)$
is the weaker condition tailored to forward projection and remains
meaningful with exact dependencies.  Under the unit-variance
normalization, define the local-isotropy coefficient by
\begin{equation}\label{eq:local-isotropy}
        \rho_s(C)=\sup_{|J|\le2s}\norm{C_{JJ}-I}_{\mathrm{op}}.
\end{equation}
It compares $C$ with $I$ and is needed only to reduce conditional
products to a static list.

The same construction gives a finite-sample certificate for the
regression search itself.  When $\mathbf x^\top\mathbf x>0$, define
\begin{equation}\label{eq:observable-forward-scores}
\begin{gathered}
        \hat{\mathbf r}=\mathbf y-\hat\beta_{0,\emptyset}\mathbf x,
        \qquad
        \hat Z_{2n}=\frac{\mathbf X^\top\hat{\mathbf r}}{\sqrt n},\\
        \hat\sigma_0^2=\frac{\mathbf x^\top\mathbf x}{n},
        \qquad
        \hat\sigma_r^2=\frac{\hat{\mathbf r}^\top\hat{\mathbf r}}{n}.
\end{gathered}
\end{equation}
Apply Theorem~\ref{thm:projection-products} with $C=\hat C$,
$a=Z_{1n}$, $b=\hat Z_{2n}$, any positive scales
$\sigma_a,\sigma_b$, and a supplied
$0<\mu\le\mu_s(\hat C)$.  Denote the resulting path by
$J_\bullet$, form $\hat z_{\mathrm a},\hat z_{\mathrm p}$ as in
\eqref{eq:score-polarization}, and put
\begin{equation}\label{eq:observable-forward-certificate}
\begin{aligned}
        \hat{\mathcal D}_n
        =\mathcal D_{\hat C,\mu}
        (\hat z_{\mathrm a},\hat z_{\mathrm p};J_\bullet),\\
        \mathcal U_{1n}
        &=\mathcal U_{\hat C,\mu}(Z_{1n};J_\bullet),
        \qquad
        \mathcal U_{2n}
        =\mathcal U_{\hat C,\mu}(\hat Z_{2n};J_\bullet).
\end{aligned}
\end{equation}
On $\hat\sigma_0\hat\sigma_r>0$, the sample standard deviations
in \eqref{eq:observable-forward-scores} give a canonical choice of the
two polarization scales.

\begin{theorem}[Finite-sample regression certificate]
\label{thm:observable-forward-search}
If $\mathcal U_{1n}<n\hat\sigma_0^2$, then every searched denominator
is positive and
\begin{equation}\label{eq:observable-forward-regret}
\begin{split}
        0&\le
        \max_{|S|=s}\hat\beta_{0,S}
        -\hat\beta_{0,J_s}\\
        &\le
        \frac{\sigma_a\sigma_b\hat{\mathcal D}_n}
        {2n\hat\sigma_0^2}
        +\frac{2\mathcal U_{1n}
        \sqrt{\mathcal U_{1n}\mathcal U_{2n}}}
        {n\hat\sigma_0^2
        (n\hat\sigma_0^2-\mathcal U_{1n})}.
\end{split}
\end{equation}
Every quantity on the right is computed from
$(\mathbf x,\mathbf y,\mathbf X)$, the chosen
scales, and the supplied design bound $\mu$.
\end{theorem}

\begin{proof}
The score vectors $Z_{1n}$ and $\hat Z_{2n}$ are compatible with
$\hat C$.  Since $\mathbf x^\top\hat{\mathbf r}=0$, partitioned
regression gives, for every support,
\begin{equation}\label{eq:exact-observable-ratio}
        \hat\beta_{0,S}-\hat\beta_{0,\emptyset}
        =
        \frac{-Z_{1n,S}^\top\hat C_{SS}^{-1}
        \hat Z_{2n,S}}
        {n\hat\sigma_0^2
        -Z_{1n,S}^\top\hat C_{SS}^{-1}Z_{1n,S}}.
\end{equation}
Equation~\eqref{eq:projection-upper-bound} bounds the two score energies
by $\mathcal U_{1n}$ and $\mathcal U_{2n}$.  Cauchy--Schwarz on the
compressed range therefore bounds the absolute numerator in
\eqref{eq:exact-observable-ratio} by
$\sqrt{\mathcal U_{1n}\mathcal U_{2n}}$.  Replacing its denominator by
$n\hat\sigma_0^2$ changes every support value by at most
\begin{equation*}
        \frac{\mathcal U_{1n}
        \sqrt{\mathcal U_{1n}\mathcal U_{2n}}}
        {n\hat\sigma_0^2
        (n\hat\sigma_0^2-\mathcal U_{1n})}.
\end{equation*}
Theorem~\ref{thm:projection-products} bounds the common-denominator
regret of $J_s$ by
$\sigma_a\sigma_b\hat{\mathcal D}_n/
(2n\hat\sigma_0^2)$.  Applying the denominator bound once to the
maximizing support and once to $J_s$ proves
\eqref{eq:observable-forward-regret}.
\end{proof}

Theorem~\ref{thm:observable-forward-search} is deterministic and retains
the support-dependent denominator, so it applies without sparse
treatment alignment, including observational and signal-dominated
designs.

The same forward path has vanishing first-order regret under a
block-independent Gaussian score law, without local isotropy.  Suppose
that, for some $\gamma_{\mathrm{sc}}\in(-1,1)$,
\begin{equation}\label{eq:kronecker-score-law}
 \begin{pmatrix}a/\sigma_a\\ b/\sigma_b\end{pmatrix}
 \sim N\!\left(0,
 \begin{pmatrix}1&\gamma_{\mathrm{sc}}\\
 \gamma_{\mathrm{sc}}&1\end{pmatrix}\otimes C\right).
\end{equation}

\begin{theorem}[Block-covariance forward search]
\label{thm:block-forward-search}
Suppose \eqref{eq:kronecker-score-law} holds and, after a permutation,
$C$ is block diagonal with $N_{\mathrm{blk}}$ positive-rank blocks.  If
\begin{equation}\label{eq:block-forward-regime}
 s=o(N_{\mathrm{blk}}),
 \qquad
 \log(p/N_{\mathrm{blk}})=o(L),
\end{equation}
then the path \eqref{eq:forward-projection-path} satisfies
\begin{equation}\label{eq:block-forward-gap}
 \norm{z_{\mathrm a}}_{C,s}^2
 -(z_{\mathrm a})_{J_s}^\top C_{J_sJ_s}^{-1}
 (z_{\mathrm a})_{J_s}
 =o_{\Prob}(sL)
\end{equation}
and
\begin{equation}\label{eq:block-forward-product}
\begin{split}
 T_C(a,b)
 &=Q_{J_s}(a,b)+o_{\Prob}(\sigma_a\sigma_b sL)\\
 &=\sigma_a\sigma_b(1-\gamma_{\mathrm{sc}})sL
   +o_{\Prob}(\sigma_a\sigma_b sL).
\end{split}
\end{equation}
The blocks may be singular and need not be well conditioned.
\end{theorem}

\begin{proof}
The polarized scores are independent with
$z_{\mathrm a}\sim N\{0,(1-\gamma_{\mathrm{sc}})C\}$ and
$z_{\mathrm p}\sim N\{0,(1+\gamma_{\mathrm{sc}})C\}$.  Put
$\widetilde z=z_{\mathrm a}/\sqrt{1-\gamma_{\mathrm{sc}}}$.
For each fixed size-$s$ support,
$\widetilde z_S^\top C_{SS}^{-1}\widetilde z_S$ is chi-square with at most
$s$ degrees of freedom.  A chi-square tail bound and a union over the
supports give
\begin{equation}\label{eq:block-forward-upper}
 \norm{\widetilde z}_{C,s}^2
 \le\{2+o_{\Prob}(1)\}sL.
\end{equation}

Choose one positive-variance coordinate from each block.  Their
standardized scores are independent standard Gaussians, so Gaussian-tail
and binomial concentration bounds give
\begin{equation}\label{eq:block-forward-lower}
 \sum_{t=1}^s
 \left\{\frac{\widetilde z_j^2}{C_{jj}}:
 j\text{ is a block representative}\right\}_{(t)}
 \ge\{2-o_{\Prob}(1)\}s\log(eN_{\mathrm{blk}}/s)
 =\{2-o_{\Prob}(1)\}sL.
\end{equation}
At step $t$, at most $t-1$ blocks have been touched.  An untouched
representative among the $t$ largest retains its marginal projection gain,
so the forward gain is at least the $t$th largest representative square.
Telescoping \eqref{eq:block-forward-lower} and comparing with
\eqref{eq:block-forward-upper} proves \eqref{eq:block-forward-gap}.

Conditionally on the active path,
\begin{equation*}
 \frac{(z_{\mathrm p})_{J_s}^\top C_{J_sJ_s}^{-1}
 (z_{\mathrm p})_{J_s}}{1+\gamma_{\mathrm{sc}}}
 \sim\chi^2_{\operatorname{rank}(C_{J_sJ_s})},
\end{equation*}
because $z_{\mathrm p}$ is independent of $z_{\mathrm a}$.  This passive
energy is $O_{\Prob}(s)=o_{\Prob}(sL)$.  The polarization identity
\eqref{eq:conditional-product-polarization} and the exact error budget in
the proof of Theorem~\ref{thm:projection-products} give the first equality
in \eqref{eq:block-forward-product}.  The selected active energy is
$2(1-\gamma_{\mathrm{sc}})sL+o_{\Prob}(sL)$, which gives the second.
\end{proof}

Condition \eqref{eq:kronecker-score-law} is an assumption on the score law,
not a consequence of block-diagonal $C$.  It holds, for example, for
Gaussian controls independent of $(x,r)$ with $d=h=0$, in which case
$\gamma_{\mathrm{sc}}=0$.

\subsection{Separable top-s reduction}

When $C=I$, the products in \eqref{eq:conditional-product} no longer
depend on $S$, so
\begin{equation}\label{eq:top-s-products}
        T_I(a,b)
        =\sum_{t=1}^s\{-a_jb_j\}_{(t)},
\end{equation}
where braces denote decreasing order statistics.  The departure from
this static product list is measured by $\rho_s(C)$ in
\eqref{eq:local-isotropy}.
\begin{theorem}[Local-isotropy reduction]\label{thm:static-products}
If $\rho_s(C)<1$, then, for every $a,b\in\R^p$,
\begin{equation}\label{eq:static-product-bound}
        |T_C(a,b)-T_I(a,b)|
        \le\rho_s(C)\norm{a}_{C,s}\norm{b}_{C,s}.
\end{equation}
\end{theorem}

\begin{proof}
For a fixed $S$, put
$\widetilde a=C_{SS}^{-1/2}a_S$ and
$\widetilde b=C_{SS}^{-1/2}b_S$.  Then
\begin{equation*}
        Q_S(a,b)+a_S^\top b_S
        =\widetilde a^\top(C_{SS}-I)\widetilde b,
\end{equation*}
whose absolute value is at most the right side of
\eqref{eq:static-product-bound}.  Taking the uniform bound over $S$ and
using $|\max_S x_S-\max_S y_S|\le\max_S|x_S-y_S|$ proves the result.
\end{proof}

\begin{theorem}[Empirical top-$s$ approximation]
\label{thm:empirical-top-s}
Suppose Assumption~\ref{ass:sparse-moments},
\eqref{eq:local-treatment}, $\delta\le\delta_0$, and
$\rho_s(C)<1$ hold.  Use the empirical scores in
\eqref{eq:empirical-scores} and the observable quantities in
\eqref{eq:observable-forward-scores}.
Then, with probability at least $1-Ce^{-cm}$,
\begin{equation}\label{eq:empirical-top-s}
\begin{split}
\bigg|\max_{|S|=s}
        (\hat\beta_{0,S}-\hat\beta_{0,\emptyset})
        -\frac1{n\hat\sigma_0^2}
        \sum_{t=1}^s
        \{-Z_{1n,j}\hat Z_{2n,j}:j\in[p]\}_{(t)}\bigg|
        \le C\{\delta+\rho_s(C)\}a_n^{\mathrm{score}}.
\end{split}
\end{equation}
Consequently the approximation error is
$o_{\Prob}(a_n^{\mathrm{score}})$ whenever
$m=o(n)$ and $\rho_s(C)=o(1)$.
\end{theorem}

\begin{proof}
The empirical part of Theorem~\ref{thm:chaos-degeneration} changes the
searched coefficient by at most $C\delta a_n^{\mathrm{score}}$ on the
sparse score event.  On the same event,
\begin{align*}
        \norm{Z_{1n}}_{C,s}&\le C\sigma_0\sqrt m,\\
        \norm{Z_{2n}}_{C,s}&\le
        C\{\sqrt n\,\widetilde\Psi_s+\sigma_r\sqrt m\}.
\end{align*}
The first bound uses
$\sqrt n\norm{d}_{C,s}\lesssim\sigma_0\sqrt m$, which follows from
\eqref{eq:local-treatment}; the second follows from
Lemma~\ref{lem:sparse-hybrid-moments}.  Theorem~\ref{thm:static-products}
therefore bounds the difference between $T_C(Z_{1n},Z_{2n})$ and
$T_I(Z_{1n},Z_{2n})$ by
$Cn\sigma_0^2\rho_s(C)a_n^{\mathrm{score}}$.
Equation~\eqref{eq:top-s-products} identifies the latter with the
top-$s$ sum formed with $Z_{2n}$.  Finally,
\begin{equation*}
        \hat Z_{2n}
        =Z_{2n}-(\hat\beta_{0,\emptyset}-\beta_\emptyset^*)Z_{1n}.
\end{equation*}
On the same event,
$|\hat\beta_{0,\emptyset}-\beta_\emptyset^*|\le
C(\sigma_r/\sigma_0)\delta$ and
$\max_{|S|=s}\norm{Z_{1n,S}}^2
\le C\{1+\rho_s(C)\}\sigma_0^2m$.  The deterministic inequality
$|\max_S a_S-\max_S b_S|\le\max_S|a_S-b_S|$ therefore bounds the
normalized change in the top-$s$ sum by
$C\{1+\rho_s(C)\}(\sigma_r/\sigma_0)\delta^3$, which is bounded by
$C\{\delta+\rho_s(C)\}a_n^{\mathrm{score}}$.
The same event has
$|\hat\sigma_0^2/\sigma_0^2-1|\le C\delta$.  Cauchy--Schwarz and
the two sparse score bounds show that the top-$s$ sum divided by
$n\sigma_0^2$ is at most
$C\{1+\rho_s(C)\}a_n^{\mathrm{score}}$ in absolute value.  Replacing
$\sigma_0^2$ by $\hat\sigma_0^2$ therefore stays within the stated
error.  The fixed-$L_q$ versions of the same sparse-moment and score
bounds give, without requiring $m\to\infty$,
\begin{equation*}
        O_{\Prob}\!\left(
        \{\delta+\rho_s(C)\}a_n^{\mathrm{score}}\right)
\end{equation*}
for the total approximation error.  This proves the final assertion.
\end{proof}

Coordinatewise score bounds give a separable special case.  Suppose
Assumption~\ref{ass:sparse-moments} holds, $m=o(n)$,
$\rho_s(C)=o(1)$, and
\begin{equation}\label{eq:coordinate-score-cap}
        \max_{j\le p}\frac{\sqrt n\,|d_j|}{\sigma_0\sqrt L}
        +\max_{j\le p}\frac{\sqrt n\,|h_j|}{\sigma_r\sqrt L}
        =O(1),
\end{equation}
then local isotropy gives
$\sqrt{\phi_{2s}}\lesssim\delta$ and
$\widetilde\Psi_s\lesssim\sigma_r\delta$.  Hence
\eqref{eq:empirical-top-s} has the clean separable form
\begin{equation}\label{eq:observable-separable}
        \max_{|S|=s}(\hat\beta_{0,S}-\hat\beta_{0,\emptyset})
        =\frac1{n\hat\sigma_0^2}
        \sum_{t=1}^s\{-Z_{1n,j}\hat Z_{2n,j}\}_{(t)}
        +o_{\Prob}\!\left(\frac{\sigma_r}{\sigma_0}\frac mn\right).
\end{equation}
The same deterministic bound and the Gaussian sparse-score bounds,
together with Theorem~\ref{thm:chaos-degeneration}, replace
the empirical maximum by
$(n\sigma_0^2)^{-1}\sum_{t\le s}\{-Z_{1j}Z_{2j}\}_{(t)}$ with
$o_{\Prob}(a_n^{\mathrm{score}})$ error whenever
$\rho_s(C)=o(1)$.  Thus local isotropy is the step that permits
approximation of the covariance-aware search by one static top-$s$ sort.

\subsection{Independent-coordinate comparison}

The Gaussian products in this top-$s$ representation need not be
independent.
Their dependence is governed by the covariance of the two score
channels, which can be separated from the covariance $C$ of the controls.
Define the standardized centered row score and its covariance by
\begin{equation}\label{eq:score-covariance}
        \chi^\circ=
        \begin{pmatrix}
        (Xx-d)/\sigma_0\\
        (Xr-h)/\sigma_r
        \end{pmatrix},
        \qquad \Omega=\Cov(\chi^\circ),
\end{equation}
and put
\begin{equation}\label{eq:score-isotropy}
        \rho_s^{\mathrm{score}}
        =\sup_{|J|\le2s}
        \norm{\Omega_{(1,2)J,(1,2)J}-I_{2|J|}}_{\mathrm{op}}.
\end{equation}
Here $(1,2)J$ includes both score coordinates associated with the
controls in $J$.

For centered variables, write
\begin{equation*}
        \kappa_4(q_1,q_2,q_3,q_4)
        =\E[q_1q_2q_3q_4]-\E[q_1q_2]\E[q_3q_4]
        -\E[q_1q_3]\E[q_2q_4]-\E[q_1q_4]\E[q_2q_3].
\end{equation*}
The sparse mean cap used below is
\begin{equation}\label{eq:score-mean-cap}
        \sup_{|J|\le s}\left\{
        \frac{n\norm{d_J}^2}{\sigma_0^2}
        +\frac{n\norm{h_J}^2}{\sigma_r^2}
        \right\}\lesssim m.
\end{equation}
Writing $\sigma_x=\sigma_0$, when $m=o(n)$, local isotropy,
\eqref{eq:score-mean-cap}, and the sparse fourth-cumulant condition
\begin{equation*}
        \sup_{|J|\le2s}
        \max_{a,b\in\{x,r\}}
        \norm{\left[
        \frac{\kappa_4(X_j,a,X_k,b)}{\sigma_a\sigma_b}
        \right]_{j,k\in J}}_{\mathrm{op}}=o(1)
\end{equation*}
imply $\rho_s^{\mathrm{score}}=o(1)$.  Local isotropy alone does not:
fourth cumulants govern the covariance of the score products, and
entrywise cumulant bounds need not control their sparse operator norms.

\begin{theorem}[Independent-coordinate Gaussian comparison]
\label{thm:gaussian-top-s-decoupling}
Suppose $m\to\infty$, $\rho_s^{\mathrm{score}}=o(1)$, and
\eqref{eq:score-mean-cap} holds.
Let $(\xi_1,\xi_2)$ have mutually independent coordinates with
\begin{equation*}
        \xi_{1j}\sim N(\sqrt n\,d_j,\sigma_0^2),
        \qquad
        \xi_{2j}\sim N(\sqrt n\,h_j,\sigma_r^2).
\end{equation*}
The two top-$s$ functionals admit a coupling for which
\begin{equation}\label{eq:gaussian-score-decoupling}
\begin{split}
        \sum_{t=1}^s\{-Z_{1j}Z_{2j}\}_{(t)}
        -\sum_{t=1}^s\{-\xi_{1j}\xi_{2j}\}_{(t)}
        =o_{\Prob}(\sigma_0\sigma_r m).
\end{split}
\end{equation}
\end{theorem}

\begin{proof}
In standardized coordinates write
\begin{equation*}
        F(z)=\max_{|S|=s}-z_{1S}^\top z_{2S}
\end{equation*}
and smooth this maximum by
\begin{equation*}
        F_\tau(z)=\frac1\tau
        \log\sum_{|S|=s}\exp\{-\tau z_{1S}^\top z_{2S}\}.
\end{equation*}
Since $\log\binom ps\le m$, the smoothing error is at most $m/\tau$.
Gaussian interpolation between the laws with covariance $I_{2p}$ and
$\Omega$, followed by integration by parts, expresses the derivative of
$\E F_\tau$ as one half of the inner product of $\Omega-I_{2p}$ with
$\E\nabla^2F_\tau$.  The supportwise quadratic Hessians are supported on
one set of size $s$.  The softmax covariance of two supportwise gradients
is supported on their union, whose size is at most $2s$.
Equation~\eqref{eq:score-isotropy} and the sparse Gaussian moment bound
therefore give, uniformly along the interpolation,
\begin{equation*}
        \left|\frac{d}{dt}\E F_\tau\right|
        \lesssim\rho_s^{\mathrm{score}}\{s+\tau m\},
\end{equation*}
where \eqref{eq:score-mean-cap} bounds the deterministic part of the
sparse score norm.  If $\rho_s^{\mathrm{score}}=0$, the two Gaussian
laws agree.  Otherwise, taking
$\tau=(\rho_s^{\mathrm{score}})^{-1/2}$ shows that the expectations of
the two unsmoothed maxima differ by $o(m)$.

The almost-everywhere gradient of $F$ is supported on a maximizing set.
On that set,
$\nabla F^\top\Omega\nabla F
\le\{1+\rho_s^{\mathrm{score}}\}\norm{\nabla F}^2$.
Gaussian Poincar\'e and the sparse moment bound therefore give variance
$O(m)$ under both laws.  Each maximum is within
$O_{\Prob}(\sqrt m)=o_{\Prob}(m)$ of its expectation.  Placing the two
score pairs independently on the same probability space and restoring
the two marginal scales proves
\eqref{eq:gaussian-score-decoupling}.
\end{proof}

\begin{corollary}[Separable approximation]
\label{cor:separable-approximation}
Suppose Assumption~\ref{ass:sparse-moments} and
\eqref{eq:score-mean-cap} hold,
$m\to\infty$, $m=o(n)$,
$\rho_s(C)=o(1)$, and $\rho_s^{\mathrm{score}}=o(1)$.
There is an enlarged probability space on which
\begin{equation}\label{eq:independent-top-s}
        \max_{|S|=s}(\hat\beta_{0,S}-\hat\beta_{0,\emptyset})
        =\frac1{n\sigma_0^2}
        \sum_{t=1}^s\{-\xi_{1j}\xi_{2j}\}_{(t)}
        +o_{\Prob}(a_n^{\mathrm{score}}).
\end{equation}
This conclusion uses score isotropy in addition to local isotropy; the
two conditions are not interchangeable.  When the control rows are Gaussian
and independent of $(x^{(i)},r^{(i)})$, the normalized observable score
channels formed from $(\mathbf x,\hat{\mathbf r})$ are, conditionally on
$(\mathbf x,\mathbf r)$, independent $N(0,C)$ vectors.  Thus, if
$m\to\infty$, $m=o(n)$, and $\rho_s(C)=o(1)$, independent standard
Gaussian vectors $\zeta_1,\zeta_2$ may be coupled so that
\begin{equation}\label{eq:gaussian-control-separable}
        \max_{|S|=s}(\hat\beta_{0,S}-\hat\beta_{0,\emptyset})
        =\frac{\sigma_r}{n\sigma_0}
        \sum_{t=1}^s\{-\zeta_{1j}\zeta_{2j}\}_{(t)}
        +o_{\Prob}\!\left(\frac{\sigma_r}{\sigma_0}\frac mn\right).
\end{equation}
The first statement follows from Theorems~\ref{thm:chaos-degeneration},
\ref{thm:static-products}, and \ref{thm:gaussian-top-s-decoupling}.
Indeed, \eqref{eq:score-mean-cap} and $\rho_s(C)=o(1)$ imply
$\sqrt{\phi_{2s}}\lesssim\delta$.
The second follows from Theorem~\ref{thm:empirical-top-s} and the
conditional Gaussian law.  In both cases, sample-to-population scale
replacement costs $O_{\Prob}(m/\sqrt n)=o_{\Prob}(m)$.
\end{corollary}

\subsection{Fixed-budget endpoint}

The growing-budget approximation operates at the $m/n$ scale.  When
$s$ is fixed, the centered extreme order statistics fluctuate on the
finer $1/n$ scale and require a moderate-deviation argument.  The
following theorem gives the corresponding calibration.

\begin{theorem}[Fixed-budget calibration]
\label{thm:fixed-budget-calibration}
Suppose $s$ is fixed, $p\to\infty$, $\beta=0$, and treatment is
independent of the controls.  Assume the row conditions of
Assumption~\ref{ass:sparse-moments},
\begin{equation}\label{eq:fixed-budget-regime}
        \log p=o(n^{1/3}),
        \qquad
        \rho_s(C)\log p=o(1),
\end{equation}
and put
\begin{equation}\label{eq:product-centering}
        b_p=\log p-\frac12\log\log p-\frac12\log(2\pi).
\end{equation}
Then
\begin{equation}\label{eq:fixed-budget-limit}
        \frac{n\sigma_0}{\sigma_r}
        \max_{|S|=s}
        (\hat\beta_{0,S}-\hat\beta_{0,\emptyset})
        -s b_p
        \ \Longrightarrow\
        \sum_{t=1}^s(-\log H_t),
        \qquad
        H_t=\mathcal E_1+\cdots+\mathcal E_t,
\end{equation}
where the $\mathcal E_t$ are independent standard exponential
variables.  For $s=1$, the limit is standard Gumbel.
\end{theorem}

\begin{proof}
Here $r=\varepsilon$, $d=h=0$, and
$\sigma_r=\sigma_\varepsilon$.  Define the standardized score
coordinates
\begin{equation*}
        A_{n,j}=\frac1{\sigma_0\sqrt n}
        \sum_{i=1}^nX_j^{(i)}x^{(i)},
        \qquad
        E_{n,j}=\frac1{\sigma_r\sqrt n}
        \sum_{i=1}^nX_j^{(i)}r^{(i)}.
\end{equation*}
For independent standard normals $G_1,G_2$, the density of
$G_1G_2$ is $K_0(|z|)/\pi$, and therefore
\begin{equation}\label{eq:gaussian-product-tail}
        \Prob(G_1G_2>z)
        =\frac{e^{-z}}{\sqrt{2\pi z}}\{1+O(z^{-1})\}.
\end{equation}
It follows that
$p\Prob(-G_1G_2>b_p+u)\to e^{-u}$.

Fix $k$ distinct coordinates.  The associated $2k$-dimensional row
score has a moment generating function on a common neighborhood of zero
and covariance
$\operatorname{diag}(C_{JJ},C_{JJ})$.
A fixed-dimensional Gaussian coupling for sums with exponential
moments\footnote{A.~Yu. Zaitsev, ``Multidimensional version of the
results of Koml\'os, Major and Tusn\'ady for vectors with finite
exponential moments,'' \emph{ESAIM: Probability and Statistics} 2
(1998), 41--108.} therefore places its normalized sum and the matched
Gaussian vector on one space with coordinate error at most
\begin{equation*}
        r_n=C_D\frac{\log n+\log p}{\sqrt n}
\end{equation*}
with failure probability $O(p^{-D})$, for every fixed $D$.  Choose
$D>k+2$ and a fixed $K$ large enough that the matched Gaussian vector
leaves the cube $[-K\sqrt{\log p},K\sqrt{\log p}]^{2k}$ with probability
$o(p^{-k})$.  On the coupling event, a Gaussian vector in this cube puts
the normalized sum in the cube with radius $2K\sqrt{\log p}$ for all
large $n$; hence the latter also leaves its enlarged cube with probability
$o(p^{-k})$.  On the Gaussian cube, changing every coordinate by at most $r_n$ changes every
product by at most
\begin{equation*}
 \epsilon_n=2K r_n\sqrt{\log p}+r_n^2=o(1)
\end{equation*}
under $\log p=o(n^{1/3})$.  For fixed $k$,
\[
        \norm{C_{JJ}-I}_{\mathrm{op}}
        \le (k-1)\max_{j\ne\ell}|C_{j\ell}|,
\]
and each off-diagonal entry is bounded by $\rho_s(C)$.  On the same cube,
the log density ratio between covariance
$\operatorname{diag}(C_{JJ},C_{JJ})$ and the identity is
$O\{\rho_s(C)\log p\}=o(1)$.  Under the identity covariance, independence
and \eqref{eq:gaussian-product-tail} show that the union of the $k$
product boundary strips has probability
\begin{equation*}
 O(\epsilon_n p^{-k})=o(p^{-k}):
\end{equation*}
one coordinate lies in a strip of probability $O(\epsilon_n p^{-1})$
and the other $k-1$ coordinates exceed their thresholds.  The density
ratio transfers this joint strip estimate to the matched Gaussian law.
Combining it with the cube restriction and the coupling proves, uniformly over distinct
$j_1,\ldots,j_k$ and bounded $u_1,\ldots,u_k$,
\begin{equation}\label{eq:fixed-budget-factorial-tail}
\begin{split}
 &\Prob\{-A_{n,j_\ell}E_{n,j_\ell}>b_p+u_\ell,\ 1\le\ell\le k\}\\
 &\hspace{22mm}
 =\{1+o(1)\}p^{-k}
   \exp\left(-\sum_{\ell=1}^ku_\ell\right).
\end{split}
\end{equation}
The mixed factorial moments obtained from
\eqref{eq:fixed-budget-factorial-tail} identify the limiting point
process as Poisson with intensity $e^{-u}\,du$; the case $k=1$ rules
out escape at the upper boundary.  Its ordered atoms are
$-\log H_1,-\log H_2,\ldots$.

It remains only to return from the scores to the coefficient.  The proof
of Theorem~\ref{thm:empirical-top-s}, now at fixed $s$, gives
\begin{equation*}
        \frac{n\sigma_0}{\sigma_r}
        \max_{|S|=s}
        (\hat\beta_{0,S}-\hat\beta_{0,\emptyset})
        =
        \sum_{t=1}^s\{-A_{n,j}E_{n,j}\}_{(t)}+o_{\Prob}(1).
\end{equation*}
Indeed, the normalized covariance and denominator errors are
$O_{\Prob}\{\rho_s(C)\log p+(\log p)^{3/2}/\sqrt n\}=o_{\Prob}(1)$.
The continuous mapping theorem for the top $s$ atoms proves
\eqref{eq:fixed-budget-limit}.
\end{proof}

\begin{lemma}[Gaussian excursion localization]
\label{lem:gaussian-excursion-localization}
Let $g_x,g_r$ be independent $N(0,C)$ vectors and let $W$ be independent
of them with $W^\top W$ distributed as a Wishart matrix with
covariance $C$ and $n-2$ degrees of freedom.  Suppose $s$ is fixed,
$p\to\infty$, $(\log p)^2=o(n)$, and $\rho_s(C)\log p=o(1)$.  Put
\begin{equation*}
 \varpi_j=-g_{x,j}g_{r,j},
 \qquad
 \mathcal I(w)=\{j:\varpi_j\ge b_p-w\},
\end{equation*}
and choose $c_n\to\infty$ so slowly that
\begin{equation}\label{eq:gaussian-localization-slack}
 c_n=o(\log p),
 \qquad
 \frac{(\log p)^2c_n}{n}=o(1).
\end{equation}
For $|S|=s$, let
\begin{equation*}
 A_S=W_S^\top W_S+g_{r,S}g_{r,S}^\top,
 \qquad
 \Theta_S=-ng_{x,S}^\top A_S^{-1}g_{r,S},
 \qquad
 \Theta_S^\circ=-g_{x,S}^\top g_{r,S}.
\end{equation*}
There is a deterministic $K_{\mathrm{loc}}<\infty$ for which
\begin{equation}\label{eq:gaussian-localized-maximizer}
\begin{gathered}
 \Prob\!\left\{J\subseteq\mathcal I(K_{\mathrm{loc}}c_n)
 \text{ for every }J\in\arg\max_{|S|=s}\Theta_S\right\}\longrightarrow1,\\
 \sup_{\substack{S\subseteq\mathcal I(K_{\mathrm{loc}}c_n)\\|S|=s}}
 |\Theta_S-\Theta_S^\circ|=o_{\Prob}(1).
\end{gathered}
\end{equation}
The $s$ largest $\varpi_j$ lie in $[b_p-c_n,b_p+c_n]$ with probability
tending to one.
\end{lemma}

\begin{proof}
The union bound and \eqref{eq:gaussian-product-tail} give
\begin{equation*}
 \Prob\{\max_j\varpi_j>b_p+c_n\}\le Ce^{-c_n}\longrightarrow0.
\end{equation*}
Let $N_n=\sum_j1\{\varpi_j>b_p-c_n\}$.  The same tail gives
$\E N_n\asymp e^{c_n}\to\infty$.  The two-coordinate Gaussian density
comparison used in \eqref{eq:fixed-budget-factorial-tail}, now on the
region $b_p-c_n$ with $c_n=o(\log p)$, gives
\begin{equation*}
 \E\{N_n(N_n-1)\}=\{1+o(1)\}(\E N_n)^2.
\end{equation*}
Thus $N_n\ge s$ with probability tending to one by the second-moment
inequality.  This proves the asserted window for the $s$ largest
products.  The marginal tail also gives
$\E|\mathcal I(w)|\le Ce^w$ uniformly for
$c_n\le w\le(\log p)/2$.  Markov's inequality, first on dyadic $w$ and
then by monotonicity, yields
\begin{equation}\label{eq:gaussian-excursion-count}
 |\mathcal I(w)|\le e^{Kw}
\end{equation}
simultaneously in this range with probability tending to one.

Condition on $g_x,g_r$.  The set $\mathcal I(w)$ is then fixed and independent
of $W$.  A fixed-support Wishart bound followed by a union over at
most $|\mathcal I(w)|^s\le e^{Ksw}$ supports gives
\begin{equation}\label{eq:localized-gaussian-gram}
 \sup_{\substack{S\subseteq\mathcal I(w)\\|S|\le s}}
 \norm{n^{-1}W_S^\top W_S-C_{SS}}_{\mathrm{op}}
 \le K\left(\sqrt{\frac wn}+\frac wn\right).
\end{equation}
Choosing the tail constant larger than $Ks$ and summing over dyadic $w$
makes \eqref{eq:gaussian-excursion-count}--\eqref{eq:localized-gaussian-gram}
simultaneous.  This is the step for which independence of $W$ from
the excursion set is essential.

Since $\max_j(|g_{x,j}|\vee|g_{r,j}|)=O_{\Prob}(\sqrt{\log p})$, inverse
perturbation on the compressed range gives, uniformly over the same sets,
\begin{equation}\label{eq:localized-gaussian-field-error}
 \sup_{\substack{S\subseteq\mathcal I(w)\\|S|=s}}
 |\Theta_S-\Theta_S^\circ|
 \le e_n(w),
 \quad
 e_n(w)=K\log p\left\{
 \rho_s(C)+\frac{\log p}{n}
 +\sqrt{\frac wn}+\frac wn\right\}.
\end{equation}
A global sparse-Wishart bound gives the error
\begin{equation*}
 e_n^\circ
 =O_{\Prob}\!\left[\log p\left\{\rho_s(C)
 +\sqrt{\frac{\log p}{n}}+\frac{\log p}{n}\right\}\right]
 =o_{\Prob}(\log p).
\end{equation*}
Let $S_0$ contain the $s$ largest $\varpi_j$ and fix any
$S^{\mathrm{loc}}\in\arg\max_{|S|=s}\Theta_S$.
Comparison of $\Theta_{S^{\mathrm{loc}}}$ with $\Theta_{S_0}$ and the product window
first gives
\begin{equation*}
 S^{\mathrm{loc}}\subseteq
 \mathcal I\{(2s-1)c_n+2e_n^\circ\}.
\end{equation*}
Whenever $S^{\mathrm{loc}}\subseteq\mathcal I(w)$, the localized error improves the
same argument to
\begin{equation*}
 S^{\mathrm{loc}}\subseteq
 \mathcal I\{(2s-1)c_n+2e_n(w)\}.
\end{equation*}
Define
\begin{equation*}
 w_0=(2s-1)c_n+2e_n^\circ,
 \qquad
 w_{k+1}=(2s-1)c_n+2e_n(w_k).
\end{equation*}
Then $w_0=o_{\Prob}(\log p)$, and
the assumptions imply $e_n(w)\le w/4$ whenever
$K_{\mathrm{loc}}c_n\le w\le(\log p)/2$, for sufficiently large fixed
$K_{\mathrm{loc}}$.  Iteration therefore reaches
$w_k\le K_{\mathrm{loc}}c_n$.  The bounds are uniform over the maximizing
support, so the same $K_{\mathrm{loc}}$ works for every maximizer.  Finally,
\begin{equation*}
 e_n(K_{\mathrm{loc}}c_n)\lesssim
 (\log p)\rho_s(C)+\frac{(\log p)^2}{n}
 +\sqrt{\frac{(\log p)^2c_n}{n}}
 +\frac{(\log p)c_n}{n}
 \longrightarrow0,
\end{equation*}
which proves \eqref{eq:gaussian-localized-maximizer}.
\end{proof}

\begin{corollary}[Gaussian-control refinement]
\label{cor:gaussian-fixed-budget}
Suppose the assumptions of Theorem~\ref{thm:fixed-budget-calibration}
hold except for the first condition in
\eqref{eq:fixed-budget-regime}.  If the control rows are Gaussian with
law $N(0,C)$ and
\begin{equation}\label{eq:gaussian-fixed-budget-regime}
        (\log p)^2=o(n),
\end{equation}
then the limit \eqref{eq:fixed-budget-limit} still holds.
\end{corollary}

\begin{proof}
Work on $\hat\sigma_0\hat\sigma_r>0$, where
\begin{equation*}
\begin{gathered}
        \hat{\mathbf r}=\mathbf r-
        \frac{\mathbf x^\top\mathbf r}{\mathbf x^\top\mathbf x}\mathbf x,
        \qquad \hat\sigma_0^2=\frac{\mathbf x^\top\mathbf x}{n},
        \qquad \hat\sigma_r^2=
        \frac{\hat{\mathbf r}^\top\hat{\mathbf r}}{n},\\
        e_x=\frac{\mathbf x}{\sqrt n\,\hat\sigma_0},
        \qquad
        e_r=\frac{\hat{\mathbf r}}{\sqrt n\,\hat\sigma_r}.
\end{gathered}
\end{equation*}
Conditional on $(\mathbf x,\mathbf r)$, the vectors $e_x,e_r$ are orthonormal.
Put $g_x=\mathbf X^\top e_x$, $g_r=\mathbf X^\top e_r$, and let
$\mathbf X_\perp$ be the projection of $\mathbf X$ onto their orthogonal
complement.  Gaussian orthogonal invariance makes $g_x,g_r$ independent
$N(0,C)$ vectors and makes $\mathbf X_\perp$
independent of them.  With
\begin{equation*}
        A_S=\mathbf X_{\perp,S}^\top\mathbf X_{\perp,S}
        +g_{r,S}g_{r,S}^\top,
\end{equation*}
partitioned regression and the Sherman--Morrison identity give exactly
\begin{equation}\label{eq:exact-gaussian-control-field}
        \frac{n\hat\sigma_0}{\hat\sigma_r}
        (\hat\beta_{0,S}-\hat\beta_{0,\emptyset})
        =-ng_{x,S}^\top A_S^{-1}g_{r,S}.
\end{equation}

Write $\varpi_j=-g_{x,j}g_{r,j}$,
$\Theta_S=-ng_{x,S}^\top A_S^{-1}g_{r,S}$, and
$\Theta_S^\circ=-g_{x,S}^\top g_{r,S}$.
Lemma~\ref{lem:gaussian-excursion-localization}
applies conditionally on $(\mathbf x,\mathbf r)$.  Its localized maximizer and field bound,
together with the support containing the $s$ largest $\varpi_j$, give
\begin{equation*}
        \max_{|S|=s}\Theta_S
        =\sum_{t=1}^s\{\varpi_j\}_{(t)}+o_{\Prob}(1).
\end{equation*}
Equation~\eqref{eq:exact-gaussian-control-field} and the Gaussian product
limit prove the result with empirical scales.  Finally,
$\sigma_0\hat\sigma_r/(\sigma_r\hat\sigma_0)=1+O_{\Prob}(n^{-1/2})$;
multiplication by the leading $O_{\Prob}(\log p)$ maximum costs $o_{\Prob}(1)$
under \eqref{eq:gaussian-fixed-budget-regime}.
\end{proof}

The cubic restriction in \eqref{eq:fixed-budget-regime} belongs to the
non-Gaussian moderate-deviation step.  Gaussian controls permit the
quadratic range in \eqref{eq:gaussian-fixed-budget-regime}.
\section{Coefficient bounds}

The moment field separates the population landscape from sampling
variation.  Recall
\begin{equation*}
        a_n=\frac1{\sigma_0}
        \{(\sigma_r\sqrt{\phi_s}+\widetilde\Psi_s)\delta
        +\sigma_r\delta^2\},
        \qquad
        S^\star\in\mathop{\arg\max}_{|S|=s}B_S(\Gamma).
\end{equation*}
For $\epsilon\ge0$, let
\begin{equation}\label{eq:near-optimal-family}
        \mathcal N_\epsilon
        =\{S:|S|=s,\quad
        B_S(\Gamma)\ge B_{S^\star}(\Gamma)-\epsilon\}.
\end{equation}
The lower bound depends on separation in the fluctuation field, not on
the cardinality of this family alone.  Put
\begin{equation}\label{eq:linear-matched-field}
        W_S=\frac1{\sqrt n}DB_S(\Gamma)[G],
        \qquad
        d_G(S,T)^2=\E(W_S-W_T)^2.
\end{equation}

\begin{theorem}[Coefficient bounds and anchored packing]
\label{thm:coefficient-bounds}
Suppose Assumption~\ref{ass:sparse-moments} holds and
$\delta\le\delta_0$.  Then, with probability at least $1-Ce^{-cm}$,
\begin{equation}\label{eq:coefficient-upper}
        \max_{|S|=s}\hat\beta_{0,S}
        -\hat\beta_{0,S^\star}
        \le Ca_n.
\end{equation}
On the same event, simultaneously for every nonempty family
$\mathcal P$ of size-$s$ supports,
\begin{equation}\label{eq:critical-empirical-lower}
\begin{split}
        \max_{|S|=s}\hat\beta_{0,S}-\hat\beta_{0,S^\star}
        \ge{}&\max_{S\in\mathcal P}
        (\mathcal T_{n,S}-\mathcal T_{n,S^\star})
        -C\frac{\sigma_r}{\sigma_0}\delta^3.
\end{split}
\end{equation}
For a distributional lower bound, suppose in addition that $m=o(n)$.
Let the fixed deterministic family
$\mathcal P\subseteq\mathcal N_\epsilon$ have cardinality $M\ge2$
and satisfy
\begin{equation}\label{eq:coefficient-packing}
        d_G(S,T)\ge\lambda,
        \qquad S\ne T,\quad S,T\in\mathcal P.
\end{equation}
Define
\begin{equation*}
        \bar\sigma_{\mathcal P}^2
        =\max_{S\in\mathcal P}
        \Var(W_S-W_{S^\star}),
        \qquad
        \omega_n=\left(\frac mn\right)^{1/6}+\delta.
\end{equation*}
Then, for $u\ge0$ and $0<\zeta\le1$, with probability at least
$1-e^{-u^2/2}-Ce^{-cm}-C\omega_n\zeta^{-3}$,
\begin{equation}\label{eq:coefficient-packing-lower}
\begin{split}
        \max_{|S|=s}\hat\beta_{0,S}
        -\hat\beta_{0,S^\star}
        \ge{}&
        c\lambda\sqrt{\log M}
        -u\bar\sigma_{\mathcal P}-\epsilon\\
        &-C\frac{\sigma_r}{\sigma_0}\delta^2-\zeta a_n .
\end{split}
\end{equation}
\end{theorem}

\begin{proof}
Equation~\eqref{eq:intrinsic-linear-row-bound} with $H=0$, Bernstein's
inequality, and a union over the supports bound the empirical linear
field by
\begin{equation*}
        C\left\{
        \frac{\widetilde\Psi_s}{\sigma_0}\delta
        +\frac{\sigma_r}{\sigma_0}\sqrt{\phi_s}\,\delta
        +\frac{\sigma_r}{\sigma_0}\delta^2
        \right\}.
\end{equation*}
Here the Bernstein linear term is absorbed by the final summand because
$\sqrt{\phi_s}\le1$ and
$\widetilde\Psi_s\le\sigma_r$.
The first-order remainder in
Theorem~\ref{thm:critical-empirical-expansion} is of the final order.
Since $B_S(\Gamma)\le B_{S^\star}(\Gamma)$, this proves
\eqref{eq:coefficient-upper}.
The second-order maximum expansion, applied at $S$ and $S^\star$, gives
\eqref{eq:critical-empirical-lower}.

For the lower bound, the uniform expansion of the matched field gives
\begin{equation*}
        \max_{|S|=s}
        |\mathcal T^G_{n,S}-B_S(\Gamma)-W_S|
        \le C\frac{\sigma_r}{\sigma_0}\delta^2
\end{equation*}
outside an event of probability $Ce^{-cm}$.  Every
$S\in\mathcal P$ has population deficit at most $\epsilon$, so the
matched anchored maximum is at least
\begin{equation*}
        \max_{S\in\mathcal P}(W_S-W_{S^\star})
        -\epsilon-C\frac{\sigma_r}{\sigma_0}\delta^2.
\end{equation*}
Sudakov minoration and Gaussian lower-tail concentration bound the first
term from below by
$c\lambda\sqrt{\log M}-u\bar\sigma_{\mathcal P}$, except with
probability $e^{-u^2/2}$.  Finally, if $d_3(A,B)\le\omega$, a smooth
cutoff at width $\zeta$ gives
\[
        \Prob(A\le t-\zeta)
        \le\Prob(B<t)+C\omega\zeta^{-3}.
\]
Apply this inequality to the anchored matched-moment approximation,
normalized by $a_n$, to obtain
\eqref{eq:coefficient-packing-lower}.
\end{proof}

Under random assignment,
$d=0$, $h=C\beta$, and
$\sigma_r^2=\beta^\top C\beta+\sigma_\varepsilon^2$.  Hence
\begin{equation}\label{eq:experimental-upper-scale}
        a_n
        =\frac1{\sigma_0}(\Psi_s\delta+\sigma_r\delta^2)
        \asymp
        \frac{\delta}{\sigma_0}
        (\Psi_s^2+\sigma_r^2\delta^2)^{1/2}.
\end{equation}
For arbitrary $d$, the additional term
$(\sigma_r/\sigma_0)\sqrt{\phi_s}\,\delta$ is generally necessary.
Sparse treatment alignment absorbs it into the quadratic term.

The score field also gives a compact lower-bound criterion in
nonexperimental designs.  Write its covariance in blocks as
$\Sigma_{ab}=\Cov(Z_a,Z_b)$ and define
\begin{equation}\label{eq:conditional-score-law}
\begin{split}
        \Sigma_{1\mid2}
        &=\Sigma_{11}-\Sigma_{12}\Sigma_{22}^{\dagger}\Sigma_{21},\\
        \bar z_1(z)
        &=\sqrt n\,d+
        \Sigma_{12}\Sigma_{22}^{\dagger}(z-\sqrt n\,h).
\end{split}
\end{equation}
This conditional law is well-defined because
$Z_2-\sqrt n\,h$ lies in the range of $\Sigma_{22}$ almost surely.
Let $\iota_S:\R^S\to\R^p$ denote coordinate embedding, put
$q_S(z)=\iota_SC_{SS}^{-1}z_S$, $q_\emptyset(z)=0$, and write
$\norm{z}_A=(z^\top Az)^{1/2}$ for $A\succeq0$.

\begin{lemma}[Anchored score approximation]
\label{lem:anchored-score-approximation}
Suppose Assumption~\ref{ass:sparse-moments},
\eqref{eq:local-treatment}, $\delta\le\delta_0$, and $m=o(n)$ hold.
Fix $S_0$ with $|S_0|\le s$, and put
\begin{equation*}
 B_{*,S_0}=\max_{|S|=s}B_S(\Gamma)-B_{S_0}(\Gamma),
 \qquad
 \omega_n=\left(\frac mn\right)^{1/6}+\delta.
\end{equation*}
Then
\begin{equation}\label{eq:anchored-score-approximation}
\begin{split}
 d_3\Bigg(
 &\frac{\max_{|S|=s}\{B_S(\hat\Gamma)-B_{S_0}(\hat\Gamma)\}
 -B_{*,S_0}}{a_n^{\mathrm{score}}},\\
 &\frac{\{T_C(Z_1,Z_2)-Q_{S_0}(Z_1,Z_2)\}/(n\sigma_0^2)
 -B_{*,S_0}}{a_n^{\mathrm{score}}}
 \Bigg)\lesssim\omega_n.
\end{split}
\end{equation}
\end{lemma}

\begin{proof}
Under \eqref{eq:local-treatment},
$a_n^{\mathrm{score}}\le a_n\le Ca_n^{\mathrm{score}}$.
The anchored conclusion of Theorem~\ref{thm:matched-moments} therefore
remains valid at the score normalization, after changing its constant.
Put
\begin{equation*}
 R_n^G=\sup_{|S|\le s}\left|
 \mathcal T^G_{n,S}-\frac{Q_S(Z_1,Z_2)}{n\sigma_0^2}\right|.
\end{equation*}
Equation~\eqref{eq:score-reduction-truncated-error} gives
\begin{equation*}
 \E\left(1\wedge\frac{R_n^G}{a_n^{\mathrm{score}}}\right)
 \lesssim\delta.
\end{equation*}
The absolute difference between the two anchored matched maxima is at
most $2R_n^G$.  Since every test function in $d_3$ is Lipschitz, the
triangle inequality proves \eqref{eq:anchored-score-approximation}.
\end{proof}

\begin{theorem}[Conditional score packing]
\label{thm:conditional-score-packing}
Fix a deterministic integer $M\ge2$ and a deterministic
$\pi_0\in[0,1]$.  Let $\mathcal P=\mathcal P(Z_2)$ be a measurable family
containing exactly $M$ size-$s$ supports, let $S_0$ be a fixed anchor with
$|S_0|\le s$, and fix $K\ge1$.  Suppose that, on a $Z_2$-measurable event
of probability at least $1-\pi_0$,
there are measurable $\lambda_n>0$ and $\Delta_n\ge0$ such that
\begin{equation}\label{eq:conditional-score-packing}
\begin{gathered}
        \min_{\substack{S,T\in\mathcal P\\S\ne T}}
        \norm{q_S(Z_2)-q_T(Z_2)}_{\Sigma_{1\mid2}}
        \ge\lambda_n,\\
        \max_{S\in\mathcal P}
        \norm{q_S(Z_2)-q_{S_0}(Z_2)}_{\Sigma_{1\mid2}}
        \le K\lambda_n,\\
        \min_{S\in\mathcal P}
        \{-\bar z_1(Z_2)^\top
        [q_S(Z_2)-q_{S_0}(Z_2)]\}\ge-\Delta_n.
\end{gathered}
\end{equation}
Then, with probability at least $1-\pi_0-Ce^{-c\log M}$,
\begin{equation}\label{eq:conditional-score-lower}
        T_C(Z_1,Z_2)-Q_{S_0}(Z_1,Z_2)
        \ge c\lambda_n\sqrt{\log M}-\Delta_n.
\end{equation}
If Assumption~\ref{ass:sparse-moments},
\eqref{eq:local-treatment}, and $\delta\le\delta_0$ also hold, the same
event, outside an additional event of probability $Ce^{-cm}$, has
\begin{equation}\label{eq:conditional-matched-coefficient-lower}
\begin{split}
        \max_{|S|=s}\mathcal T^G_{n,S}-\mathcal T^G_{n,S_0}
        \ge{}&\frac{c\lambda_n\sqrt{\log M}-\Delta_n}
        {n\sigma_0^2}\\
        &-C\left\{
        \frac{\widetilde\Psi_s}{\sigma_0}\delta^2
        +\frac{\sigma_r}{\sigma_0}\delta^3\right\}.
\end{split}
\end{equation}
Under these same additional moment conditions, suppose moreover that
$m=o(n)$ and that there are deterministic sequences
$\underline\lambda_n>0$ and $\overline\Delta_n\ge0$ for which
$\lambda_n\ge\underline\lambda_n$ and
$\Delta_n\le\overline\Delta_n$ on the packing event.  For
$0<\zeta\le1$, with probability at least
\begin{equation*}
        1-\pi_0-Ce^{-c\log M}
        -C\left\{(m/n)^{1/6}+\delta\right\}\zeta^{-3},
\end{equation*}
\begin{equation}\label{eq:conditional-empirical-coefficient-lower}
\begin{split}
        \max_{|S|=s}B_S(\hat\Gamma)-B_{S_0}(\hat\Gamma)
        \ge{}&\frac{c\underline\lambda_n\sqrt{\log M}
        -\overline\Delta_n}{n\sigma_0^2}
        -\zeta a_n^{\mathrm{score}}.
\end{split}
\end{equation}
\end{theorem}

\begin{proof}
Conditional on $Z_2=z$,
$Z_1\sim N\{\bar z_1(z),\Sigma_{1\mid2}\}$.  The centered
canonical distance between two support values is therefore
$\norm{q_S(z)-q_T(z)}_{\Sigma_{1\mid2}}$.  Conditional Sudakov
minoration and Gaussian lower-tail concentration, using the radius in
\eqref{eq:conditional-score-packing}, give
$c\lambda_n\sqrt{\log M}$; the final line of that display accounts for
the conditional mean.  This proves \eqref{eq:conditional-score-lower}.
The uniform score reduction proves
\eqref{eq:conditional-matched-coefficient-lower}.  On the packing event,
the deterministic envelopes give
\begin{equation*}
 \frac{T_C(Z_1,Z_2)-Q_{S_0}(Z_1,Z_2)}{n\sigma_0^2}
 \ge\ell_n,
 \qquad
 \ell_n=\frac{c\underline\lambda_n\sqrt{\log M}
 -\overline\Delta_n}{n\sigma_0^2},
\end{equation*}
except with conditional probability $Ce^{-c\log M}$.  Apply the smooth
cutoff inequality to Lemma~\ref{lem:anchored-score-approximation} at the
deterministic threshold
$(\ell_n-B_{*,S_0})/a_n^{\mathrm{score}}$, where
$B_{*,S_0}=\max_{|S|=s}B_S(\Gamma)-B_{S_0}(\Gamma)$.  The cutoff has width $\zeta$,
so its first three derivatives are bounded by $C\zeta^{-3}$.  This proves
\eqref{eq:conditional-empirical-coefficient-lower} with the stated
probability.
\end{proof}

For example, if $\Sigma_{1\mid2}\asymp\sigma_0^2C$ on $2s$-sparse
vectors, a pool of $N$ outcome scores of size
$b_n\asymp\sigma_r\sqrt{\log(eN/s)}$ gives
$\lambda_n\asymp\sigma_0b_n\sqrt s$.  With
$\log(eN/s)\asymp L$ and negligible $\Delta_n$, the resulting coefficient
lower bound is $(\sigma_r/\sigma_0)m/n$.  Under random assignment the
covariance comparison and mean-retention condition hold automatically;
in an observational design they are the necessary score-overlap and
landscape-retention conditions.

At the critical boundary, a direct conditional argument does not require
$m=o(n)$.  A sequence of centered treatment vectors
$\mathbf x\in\R^n$ with covariance $\sigma_0^2I_n$ has uniform exponential
multiplier minoration if, for each fixed $K$, there are constants
$c_K,C_K>0$ such that every finite family $a_1,\ldots,a_M$ satisfying
\begin{equation}\label{eq:multiplier-minoration-geometry}
        \min_{k\ne\ell}\sigma_0\norm{a_k-a_\ell}_2\ge\lambda,
        \qquad
        \max_k\sigma_0\norm{a_k}_2\le K\lambda
\end{equation}
obeys
\begin{equation}\label{eq:multiplier-minoration}
        \Prob\left(
        \max_{k\le M}a_k^\top\mathbf x
        \ge c_K\lambda\sqrt{\log M}\right)
        \ge1-C_Ke^{-c_K\log M}.
\end{equation}
Gaussian treatment satisfies this property.  General sub-Gaussian
treatment need not.

Define the embedded empirical outcome score
\begin{equation}\label{eq:embedded-outcome-score}
        q_{n,S}=\iota_SC_{SS}^{-1}Z_{2n,S},
        \qquad |S|=s.
\end{equation}

\begin{theorem}[Experimental matching from score packing]
\label{thm:experimental-score-matching}
Suppose Assumption~\ref{ass:sparse-moments} holds and
$\delta\le\delta_0$.  Suppose $\mathbf x$ is independent of
$\mathcal F=\sigma(\mathbf X,\mathbf r)$ and has uniform exponential multiplier
minoration.  Fix a deterministic integer $M\ge2$ and a deterministic
$\pi_0\in[0,1]$.  Let $\mathcal P$ be an $\mathcal F$-measurable family
containing exactly $M$ size-$s$ supports, let $K\ge1$ be fixed, and let
$\Lambda_n>0$ be $\mathcal F$-measurable.  Suppose that, on an
$\mathcal F$-measurable event of probability at least $1-\pi_0$,
\begin{equation}\label{eq:experimental-score-packing}
        \min_{\substack{S,T\in\mathcal P\\S\ne T}}
        \norm{q_{n,S}-q_{n,T}}_C\ge\Lambda_n,
        \qquad
        \max_{S\in\mathcal P}\norm{q_{n,S}}_C\le K\Lambda_n .
\end{equation}
Then, with probability at least
$1-\pi_0-Ce^{-c\log M}-Ce^{-cm}$,
\begin{equation}\label{eq:experimental-packing-lower}
\begin{split}
        \max_{|S|=s}
        (\hat\beta_{0,S}-\hat\beta_{0,\emptyset})
        \ge{}&
        c\frac{\Lambda_n\sqrt{\log M}}{n\sigma_0}\\
        &-C\left\{
        \frac{\Psi_s}{\sigma_0}\delta^2
        +\frac{\sigma_r}{\sigma_0}\delta^3
        \right\}.
\end{split}
\end{equation}
If, in addition,
\begin{equation}\label{eq:experimental-matching-bridge}
        \log M\ge cm,
        \qquad
        \Lambda_n\ge
        c\sqrt n(\Psi_s^2+\sigma_r^2\delta^2)^{1/2},
\end{equation}
and $\pi_0\le Ce^{-cm}$, then, for sufficiently small $\delta_0$,
\begin{equation}\label{eq:experimental-matching}
        \max_{|S|=s}
        (\hat\beta_{0,S}-\hat\beta_{0,\emptyset})
        \asymp
        \frac{\delta}{\sigma_0}
        (\Psi_s^2+\sigma_r^2\delta^2)^{1/2}
\end{equation}
with probability at least $1-Ce^{-cm}$.
\end{theorem}

\begin{proof}
The empirical score reduction gives, uniformly over the supports,
\begin{equation*}
        B_S(\hat\Gamma)
        =-\frac{Z_{1n,S}^\top C_{SS}^{-1}Z_{2n,S}}
        {n\sigma_0^2}+R_{n,S},
        \qquad
        \max_S|R_{n,S}|
        \le C\left\{
        \frac{\Psi_s}{\sigma_0}\delta^2
        +\frac{\sigma_r}{\sigma_0}\delta^3\right\}
\end{equation*}
outside an event of probability $Ce^{-cm}$.  Conditional on
$\mathcal F$, the leading field is centered and linear in $\mathbf x$:
\begin{equation*}
        -\frac{Z_{1n,S}^\top C_{SS}^{-1}Z_{2n,S}}
        {n\sigma_0^2}
        =-\frac{\mathbf x^\top\mathbf Xq_{n,S}}
        {n^{3/2}\sigma_0^2}.
\end{equation*}
Its conditional distance is
\[
        \frac{\norm{q_{n,S}-q_{n,T}}_{\hat C}}{n\sigma_0},
        \qquad \hat C=\mathbf X^\top\mathbf X/n.
\]
Relative restricted isometry makes this equivalent to the $C$-norm
on every $2s$-sparse difference.  Multiplier minoration and
\eqref{eq:experimental-score-packing} therefore yield the first term in
\eqref{eq:experimental-packing-lower}; the score-reduction remainder
gives the second.  Under \eqref{eq:experimental-matching-bridge}, the
remainder is smaller by a factor $O(\delta)$.  The matching upper bound
is \eqref{eq:experimental-upper-scale}.
\end{proof}

The packing premise has the same top-$s$ construction in both the noise
and signal regimes.  A greedy packing of the constant-weight Hamming
slice supplies the deterministic code used below.

\begin{lemma}[Measurable top-$s$ score pool]
\label{lem:top-s-score-pool}
Fix a deterministic $N\ge2s$ and a deterministic constant-weight code
$\mathcal A_N\subseteq\{A\subseteq[N]:|A|=s\}$ satisfying
\begin{equation*}
 \log|\mathcal A_N|\ge cs\log(eN/s),
 \qquad |A\mathbin\triangle B|\ge cs,
 \quad A\ne B.
\end{equation*}
Let $\mathcal J$ be a measurable random subset of $[p]$ with
$|\mathcal J|\ge N$, write $j_1<\cdots<j_N$ for its first $N$ elements,
and define
\begin{equation*}
 \mathcal P(\mathcal J)
 =\bigl\{\{j_k:k\in A\}:A\in\mathcal A_N\bigr\}.
\end{equation*}
Then $\mathcal P(\mathcal J)$ is measurable and has deterministic
cardinality $M=|\mathcal A_N|$.  Suppose, for $z=Z_{2n}$,
\begin{equation*}
 b_n\le|z_j|\le K_{\mathrm{pool}}b_n,
 \qquad j\in\mathcal J,
 \qquad
 \rho_s(C)\le c_{\mathrm{pool}}/K_{\mathrm{pool}},
\end{equation*}
where $K_{\mathrm{pool}}$ is fixed and $c_{\mathrm{pool}}>0$ is a
sufficiently small fixed constant.  Then, with $L_N=\log(eN/s)$,
\begin{equation}\label{eq:top-s-score-pool}
\begin{gathered}
 \log M\ge csL_N,\qquad
 \min_{\substack{S,T\in\mathcal P(\mathcal J)\\S\ne T}}
 \norm{q_{n,S}-q_{n,T}}_C\ge cb_n\sqrt s,\\
 \max_{S\in\mathcal P(\mathcal J)}
 \norm{q_{n,S}}_C\le Cb_n\sqrt s.
\end{gathered}
\end{equation}
\end{lemma}

\begin{proof}
Such a deterministic code exists.  Indeed, greedily pack the
constant-weight slice and delete after each choice every set at symmetric
difference less than $c_0s$.  Standard binomial bounds give, for a
sufficiently small absolute $c_0>0$,
\begin{equation*}
 \frac{\binom Ns}
 {\sum_{k\le c_0s/2}\binom sk\binom{N-s}{k}}
 \ge \exp\{cs\log(eN/s)\}.
\end{equation*}
Put $\widetilde q_S=\iota_Sz_S$ and $\rho=\rho_s(C)$.  Local isotropy
and the inverse identity give
\begin{equation*}
 \norm{q_{n,S}-\widetilde q_S}_2
 \le\frac{\rho}{1-\rho}\norm{z_S}_2
 \le2\rho K_{\mathrm{pool}}b_n\sqrt s.
\end{equation*}
For distinct codewords,
\begin{equation*}
 \norm{\widetilde q_S-\widetilde q_T}_2^2
 =\sum_{j\in S\mathbin\triangle T}z_j^2
 \ge cb_n^2s.
\end{equation*}
Choosing $c_{\mathrm{pool}}$ sufficiently small therefore gives
$\norm{q_{n,S}-q_{n,T}}_2\ge cb_n\sqrt s$.  The difference is supported
on at most $2s$ coordinates, so local isotropy converts this Euclidean
bound into the $C$-norm bound in \eqref{eq:top-s-score-pool}.  Finally,
\begin{equation*}
 \norm{q_{n,S}}_C^2=z_S^\top C_{SS}^{-1}z_S
 \le\frac{\norm{z_S}_2^2}{1-\rho}
 \le Cb_n^2s.
\end{equation*}
\end{proof}

Thus \eqref{eq:experimental-matching} follows on this pool event when
$L_N\gtrsim L$ and
\begin{equation*}
 b_n\sqrt s\gtrsim
 \sqrt n(\Psi_s^2+\sigma_r^2\delta^2)^{1/2}.
\end{equation*}
The noise and signal choices are respectively
$b_n\asymp\sigma_r\sqrt L$ and
$b_n\asymp\sqrt n\,\Psi_s/\sqrt s$.

\section{Studentized search}

The linear specialization \eqref{eq:model} is in force throughout this
section.

Let
\begin{equation*}
        v_y^{(i)}
        =\begin{pmatrix}x^{(i)}\\y^{(i)}\\X^{(i)}\end{pmatrix},
        \qquad
        \Gamma_y=\E[v_y^{(i)}(v_y^{(i)})^\top],
        \qquad
        \hat\Gamma_y=\frac1n\sum_{i=1}^n
        v_y^{(i)}(v_y^{(i)})^\top .
\end{equation*}
For a second-moment matrix indexed by $(x,y,X)$, define
\begin{equation}\label{eq:t-functional}
        \tau_S(A)
        =\frac{A_{xy\mathbin{\cdot}S}}
        {\{A_{xx\mathbin{\cdot}S}A_{yy\mathbin{\cdot}S}\}^{1/2}} .
\end{equation}
For every $|S|\le s$, define the scaled partial-correlation statistic
and its conventional least-squares transform by
\begin{align}
        \hat t_{0,S}
        &=\sqrt{n-|S|-1}\,\tau_S(\hat\Gamma_y),
        \label{eq:studentized-moment-identity}\\
        \hat t^{\mathrm{LS}}_{0,S}
        &=\sqrt{n-\operatorname{rank}(\mathbf X_S)-1}\,
        \frac{\tau_S(\hat\Gamma_y)}
        {\{1-\tau_S(\hat\Gamma_y)^2\}^{1/2}} .
        \label{eq:conventional-t-definition}
\end{align}
Set either empirical statistic to zero when its displayed denominator
vanishes.  The variance floors and sparse concentration make this convention
irrelevant on the events used below.
Assume throughout that
\begin{equation}\label{eq:t-floors}
        \inf_{|S|\le s}(\Gamma_y)_{xx\mathbin{\cdot}S}
        \ge c_x\Var(x),
        \qquad
        \inf_{|S|\le s}(\Gamma_y)_{yy\mathbin{\cdot}S}
        \ge c_y\Var(y).
\end{equation}
Write $|H|_{y,s}$ for the sparse norm in
\eqref{eq:sparse-moment-size} with $r$ replaced by $y$ and with
$|H_{yy}|/\Var(y)$ included.
The support-specific net proof of
Lemma~\ref{lem:sparse-hybrid-moments} applies to this full fixed-dimensional
target block and gives the same tail and fixed-moment bounds for
$|\hat\Gamma_y-\Gamma_y|_{y,s}$ and its deterministic hybrids.

Under the null $\beta_0=0$, put
\begin{equation*}
        \theta_s=\frac{\Psi_s^2}{\Var(y)}
\end{equation*}
and impose the local testing condition
\begin{equation}\label{eq:local-testing-condition}
        \max_{|S|\le s}|\tau_S(\Gamma_y)|
        \le K_\tau n^{-1/2}.
\end{equation}
Let
\begin{equation*}
        H_n=\hat\Gamma_y-\Gamma_y,
        \qquad
        t_{0,S}^\star
        =\sqrt{n-|S|-1}\,\tau_S(\Gamma_y),
\end{equation*}
and define the second-order field
\begin{equation}\label{eq:studentized-second-order-field}
        \mathcal T^t_{n,S}
        =t_{0,S}^\star+\sqrt{n-|S|-1}\left\{
        D\tau_S(\Gamma_y)[H_n]
        +\frac12D^2\tau_S(\Gamma_y)[H_n,H_n]
        \right\}.
\end{equation}

\begin{theorem}[Studentized moment reduction]
\label{thm:studentized-search}
Suppose the relative sub-Gaussian condition of
Assumption~\ref{ass:sparse-moments} holds for $v_y$, the floors in
\eqref{eq:t-floors} hold, $\delta\le\delta_0$, $\beta_0=0$, and
\eqref{eq:local-testing-condition} holds.  With probability at least $1-Ce^{-cm}$,
simultaneously over $|S|\le s$,
\begin{equation}\label{eq:critical-studentized-expansion}
        \left|
        \hat t_{0,S}-\mathcal T^t_{n,S}
        \right|
        \le C\sqrt m\,\delta^2.
\end{equation}
Consequently,
\begin{equation}\label{eq:general-studentized-upper-scale}
        \max_{|S|=s}\hat t_{0,S}-\hat t_{0,\emptyset}
        \le C\left[
        \sqrt m\{\sqrt{\theta_s}+\sqrt{\phi_s}+\delta\}+1
        \right].
\end{equation}
Under sparse treatment alignment
$\sqrt{\phi_{2s}}\lesssim\delta$, this becomes
\begin{equation}\label{eq:studentized-upper-scale}
        \max_{|S|=s}\hat t_{0,S}-\hat t_{0,\emptyset}
        \le C\left[
        \sqrt m\{\sqrt{\theta_s}+\delta\}+1
        \right].
\end{equation}

The same event gives an anchored lower reduction.  For every anchor
$S_0$ with $|S_0|\le s$ and every nonempty family
$\mathcal P\subseteq\{S:|S|=s\}$,
\begin{equation}\label{eq:critical-studentized-lower}
\begin{split}
        \max_{|S|=s}\hat t_{0,S}-\hat t_{0,S_0}
        \ge{}&\max_{S\in\mathcal P}
        (\mathcal T^t_{n,S}-\mathcal T^t_{n,S_0})\\
        &-C\sqrt m\,\delta^2.
\end{split}
\end{equation}
Finally,
\begin{equation}\label{eq:least-squares-transfer}
        \max_{|S|\le s}
        |\hat t^{\mathrm{LS}}_{0,S}-\hat t_{0,S}|
        \le C\sqrt m\,\delta^2.
\end{equation}
The fixed-moment version is
\begin{equation}\label{eq:least-squares-truncated-transfer}
 \E\left[1\wedge
 \frac{\max_{|S|\le s}|\hat t^{\mathrm{LS}}_{0,S}-\hat t_{0,S}|}
 {\sqrt m\,\delta}\right]
 \le C\delta.
\end{equation}
Hence the upper and anchored lower bounds remain valid for
$\hat t^{\mathrm{LS}}$, after changing the remainder constant.
\end{theorem}

\begin{proof}
Support compression reduces every singular population block to its
range, on which the Schur complements are ordinary smooth functions.
On $|H_n|_{y,s}\le C\delta$, differentiation of
\eqref{eq:t-functional} and the floors in \eqref{eq:t-floors} give a
sharper anchored derivative bound through
Lemma~\ref{lem:residual-field-row-bounds}.  Indeed, take
$U=(x,y)^\top$, $D_U=\operatorname{diag}\{\sigma_0,\sqrt{\Var(y)}\}$,
$\varphi(R)=R_{12}/(R_{11}R_{22})^{1/2}$.  Then
\begin{equation*}
 \lambda_s\le\sqrt{\phi_s}+\sqrt{\theta_s}.
\end{equation*}
The standardized derivatives in
\eqref{eq:standardized-target-derivatives} are bounded by a constant.
Consequently,
\begin{equation}\label{eq:studentized-derivative-envelope}
\begin{split}
        \max_{|S|=s}\sup_{|H|_{y,s}\le1}
        \left|D\{\tau_S-c_{s,n}\tau_\emptyset\}
        (\Gamma_y)[H]\right|
        \lesssim\sqrt{\theta_s}+\sqrt{\phi_s}+\frac{s}{n},
        \qquad
        c_{s,n}=\sqrt{\frac{n-1}{n-s-1}}.
\end{split}
\end{equation}
Here $|1-c_{s,n}|\lesssim s/n$; the first two terms are the outcome and
treatment projection sizes.  The second and third derivatives are
uniformly bounded on the same neighborhood.  Sparse-moment
concentration puts $|H_n|_{y,s}\le C\delta$ outside an event of
probability $Ce^{-cm}$.  Relative inverse perturbation then keeps both
residual variances along $\Gamma_y+tH_n$, $0\le t\le1$, above one half of
the floors in \eqref{eq:t-floors}.  Taylor's formula with the bounded
third derivative therefore gives
\eqref{eq:critical-studentized-expansion}.  The linear and quadratic
terms in the anchored field are bounded by
\begin{equation*}
 C\sqrt m\left\{\sqrt{\theta_s}+\sqrt{\phi_s}
 +\frac{s}{n}+\delta\right\}
 \le C\sqrt m\{\sqrt{\theta_s}+\sqrt{\phi_s}+\delta\},
\end{equation*}
while the
population difference is $O(1)$ by
\eqref{eq:local-testing-condition}.  This proves
\eqref{eq:general-studentized-upper-scale}; sparse treatment alignment
absorbs $\sqrt{\phi_s}$.  Subtracting two uniform expansions and
maximizing over $\mathcal P$ proves the anchored statements.

On the same event,
$\max_{|S|\le s}|\tau_S(\hat\Gamma_y)|\le C\delta$.  If
$r_S=\operatorname{rank}(\mathbf X_S)$, then
\begin{equation*}
 \left|\sqrt{n-r_S-1}-\sqrt{n-|S|-1}\right|
 \le\frac{C(|S|-r_S)}{\sqrt n},
 \qquad
 |(1-u^2)^{-1/2}-1|\le Cu^2.
\end{equation*}
Therefore
\begin{equation*}
 \max_{|S|\le s}|\hat t^{\mathrm{LS}}_{0,S}-\hat t_{0,S}|
 \le C\left\{\frac{s\delta}{\sqrt n}+\sqrt n\,\delta^3\right\}
 \le C\sqrt m\,\delta^2,
\end{equation*}
which proves \eqref{eq:least-squares-transfer}.
The same calculation with the random sparse norm in place of its
high-probability bound, followed by the fixed-moment and exponential-tail
forms of Lemma~\ref{lem:sparse-hybrid-moments}, gives
\eqref{eq:least-squares-truncated-transfer}.  The convention on undefined
statistics and truncation at one absorb the exceptional localization event.
\end{proof}

Let $G_y$ be the centered symmetric Gaussian matrix matched to the row
moment $v_y^{(i)}(v_y^{(i)})^\top$:
\begin{equation}\label{eq:studentized-matched-moment}
 \Cov\{(G_y)_{ab},(G_y)_{cd}\}
 =\Cov\{(v_y^{(i)})_a(v_y^{(i)})_b,
 (v_y^{(i)})_c(v_y^{(i)})_d\}.
\end{equation}
For $|S|\le s$, define the matched counterpart of
\eqref{eq:studentized-second-order-field} by
\begin{equation}\label{eq:studentized-matched-field}
\begin{split}
 \mathcal T^{t,G}_{n,S}
 ={}&t_{0,S}^\star
 +\sqrt{\frac{n-|S|-1}{n}}\,
 D\tau_S(\Gamma_y)[G_y]\\
 &+\frac{\sqrt{n-|S|-1}}{2n}
 D^2\tau_S(\Gamma_y)[G_y,G_y].
\end{split}
\end{equation}
Put
\begin{equation}\label{eq:studentized-matched-scale}
 a_n^t=\sqrt m\{\sqrt{\theta_s}+\sqrt{\phi_s}+\delta\},
 \qquad
 t_{*,n}^\Delta
 =\max_{|S|=s}\{t_{0,S}^\star-t_{0,\emptyset}^\star\}.
\end{equation}

\begin{corollary}[Studentized matched approximation]
\label{cor:studentized-matched-approximation}
Under the assumptions of Theorem~\ref{thm:studentized-search}, if
$m=o(n)$, then
\begin{equation}\label{eq:studentized-matched-approximation}
 d_3\!\left(
 \frac{\max_{|S|=s}\{\hat t_{0,S}-\hat t_{0,\emptyset}\}
       -t_{*,n}^\Delta}{a_n^t},
 \frac{\max_{|S|=s}\{\mathcal T^{t,G}_{n,S}
       -\mathcal T^{t,G}_{n,\emptyset}\}-t_{*,n}^\Delta}{a_n^t}
 \right)
 \lesssim \left(\frac mn\right)^{1/6}+\delta.
\end{equation}
The same bound holds with $\hat t_{0,S}$ replaced by
$\hat t^{\mathrm{LS}}_{0,S}$ in the first argument.  Under sparse treatment
alignment, both conclusions remain valid after replacing $a_n^t$ by
\begin{equation*}
 \bar a_n^t=\sqrt m\{\sqrt{\theta_s}+\delta\}.
\end{equation*}
\end{corollary}

\begin{proof}
Apply Lemma~\ref{lem:residual-field-row-bounds} with target $(x,y)$ and
the normalized anchored functional
\begin{equation*}
 \Phi_{n,S}
 =\sqrt{\frac{n-s-1}{n}}\,\tau_S
 -\sqrt{\frac{n-1}{n}}\,\tau_\emptyset.
\end{equation*}
The specialization in the proof of
\eqref{eq:studentized-derivative-envelope} has
\begin{equation*}
 \lambda_s\le\sqrt{\theta_s}+\sqrt{\phi_s},
 \qquad
 \left|\sqrt{\frac{n-s-1}{n}}
 -\sqrt{\frac{n-1}{n}}\right|\lesssim\frac{s}{n},
 \qquad \kappa_\varphi\lesssim1.
\end{equation*}
Thus, with
$\alpha_{t,s}=\sqrt{\theta_s}+\sqrt{\phi_s}+s/n$, the lemma gives
\begin{align*}
 \sup_{|S|=s}
 \norm{D\Phi_{n,S}(\Gamma_y)[Y]
 +D^2\Phi_{n,S}(\Gamma_y)[H,Y]}_{\psi_1\mid H}
 &\lesssim\alpha_{t,s}+|H|_{y,s},\\
 \sup_{|S|=s}\norm{D^2\Phi_{n,S}(\Gamma_y)[Y,Y]}_{L_q}
 &\lesssim q(s+q).
\end{align*}
The same diagonal estimate holds for the matched Gaussian row, and the
two diagonal expectations agree.  These are precisely the three inputs
used in the diagonal split and rowwise replacement in the proof of
Theorem~\ref{thm:matched-moments}.  The natural scale before restoring
the factor $\sqrt n$ is
\begin{equation*}
 \alpha_{t,s}\delta+\delta^2
 \lesssim\delta\{\sqrt{\theta_s}+\sqrt{\phi_s}+\delta\}
 =\frac{a_n^t}{\sqrt n}.
\end{equation*}
Repeating that diagonal and off-diagonal calculation therefore gives
\eqref{eq:studentized-matched-approximation}; the cubic Taylor remainder,
$O(\sqrt n\,\delta^3)$, costs $O(\delta)$ after normalization.  Sparse treatment
alignment absorbs $\sqrt{\phi_s}$ and gives the claim with $\bar a_n^t$.
Finally, \eqref{eq:least-squares-truncated-transfer} changes either
normalized first argument by $O(\delta)$ in the $d_3$ metric, because
$a_n^t,\bar a_n^t\ge\sqrt m\,\delta$.
\end{proof}

For every $|S|\le s$, put
\begin{equation}\label{eq:studentized-directions}
\begin{gathered}
        u_{n,S}=\begin{cases}
        M_S\mathbf y/\norm{M_S\mathbf y}_2,
        &\norm{M_S\mathbf y}_2>0,\\
        0,&\norm{M_S\mathbf y}_2=0,
        \end{cases}
        \qquad
        \nu_{n,S}
        =\sqrt{\frac{n-|S|-1}{n-\operatorname{rank}(\mathbf X_S)}},\\
        u_{n,S}^{\Delta}
        =\nu_{n,S}u_{n,S}
        -\nu_{n,\emptyset}u_{n,\emptyset}.
\end{gathered}
\end{equation}

\begin{corollary}[Studentized score packing]
\label{cor:studentized-score-packing}
In addition to the assumptions of Theorem~\ref{thm:studentized-search},
suppose that $\mathbf x$ is independent of
$\mathcal F=\sigma(\mathbf X,\mathbf y)$ and has
uniform exponential multiplier minoration.  Fix a deterministic integer
$M\ge2$ and a deterministic $\pi_0\in[0,1]$.  Let
$\mathcal P$ be an $\mathcal F$-measurable family containing exactly $M$
size-$s$ supports, let $K\ge1$ be fixed, and let $\lambda_n>0$ be
$\mathcal F$-measurable.  Suppose that, on an $\mathcal F$-measurable
event of probability at least $1-\pi_0$,
\begin{equation}\label{eq:studentized-packing-event}
\begin{gathered}
        \min_{S\in\mathcal P\cup\{\emptyset\}}
        \norm{M_S\mathbf y}_2>0,\\
        \min_{\substack{S,T\in\mathcal P\\S\ne T}}
        \norm{u_{n,S}^{\Delta}-u_{n,T}^{\Delta}}_2
        \ge\lambda_n,\\
        \max_{S\in\mathcal P}\norm{u_{n,S}^{\Delta}}_2
        \le K\lambda_n .
\end{gathered}
\end{equation}
Then, with probability at least
$1-\pi_0-Ce^{-c\log M}-Ce^{-cm}$,
\begin{equation}\label{eq:studentized-packing-lower}
        \max_{|S|=s}\hat t_{0,S}-\hat t_{0,\emptyset}
        \ge c\lambda_n\sqrt{\log M}-C\sqrt m\,\delta^2 .
\end{equation}
The same bound holds with $\hat t^{\mathrm{LS}}$ in place of $\hat t$.
If, in addition,
\begin{equation}\label{eq:studentized-matching-packing}
        \log M\ge cm,
        \qquad
        \lambda_n\ge c(\theta_s+\delta^2)^{1/2},
\end{equation}
with constants for which the packing term dominates the remainder, then
\begin{equation}\label{eq:studentized-lower}
        \max_{|S|=s}\hat t_{0,S}-\hat t_{0,\emptyset}
        \ge c\sqrt m\,(\theta_s+\delta^2)^{1/2}.
\end{equation}
\end{corollary}

\begin{proof}
Let $\hat\sigma_0^2=\mathbf x^\top\mathbf x/n$ and set
\begin{equation*}
        \Delta^x_{n,S}
        =\frac{\mathbf x^\top M_S\mathbf x/
        \{n-\operatorname{rank}(\mathbf X_S)\}}
        {\mathbf x^\top\mathbf x/n}-1 .
\end{equation*}
For every support with
$\mathbf x^\top M_S\mathbf x>0$ and
$\norm{M_S\mathbf y}_2>0$, the partial-correlation identity gives exactly
\begin{equation}\label{eq:exact-studentized-direction}
        \hat t_{0,S}
        =\nu_{n,S}\frac{\mathbf x^\top u_{n,S}}{\hat\sigma_0}
        (1+\Delta^x_{n,S})^{-1/2}.
\end{equation}
Conditional quadratic and linear-form concentration, followed by a union
over the supports, gives
\begin{equation*}
        \hat\sigma_0\asymp\sigma_0,
        \qquad
        \max_{|S|=s}|\Delta^x_{n,S}|\le C\delta^2,
        \qquad
        \max_{|S|=s}|\mathbf x^\top u_{n,S}|\le C\sigma_0\sqrt m
\end{equation*}
outside an event of probability $Ce^{-cm}$.  Since
$P_S=I-M_S$ has rank at most $s$, the quadratic assertion follows
explicitly from
\begin{equation*}
 \max_{|S|=s}\mathbf x^\top P_S\mathbf x\le C\sigma_0^2m,
 \qquad
 \mathbf x^\top\mathbf x\ge c n\sigma_0^2,
\end{equation*}
because
\begin{equation*}
 1+\Delta^x_{n,S}
 =\frac{n}{n-r_S}\left(1-
 \frac{\mathbf x^\top P_S\mathbf x}
 {\mathbf x^\top\mathbf x}\right),
 \qquad r_S=\operatorname{rank}(\mathbf X_S).
\end{equation*}
The linear assertion is the corresponding union bound over the unit
vectors $u_{n,S}$.  Since
$\Delta^x_{n,\emptyset}=0$, \eqref{eq:exact-studentized-direction} yields
\begin{equation*}
        \max_{|S|=s}\left|
        \hat t_{0,S}-\hat t_{0,\emptyset}
        -\frac{\mathbf x^\top u_{n,S}^{\Delta}}{\hat\sigma_0}
        \right|
        \le C\sqrt m\,\delta^2 .
\end{equation*}
Conditionally on $\mathcal F$, the process
$\mathbf x^\top u_{n,S}^{\Delta}/\sigma_0$ is a multiplier process with canonical
distance $\norm{u_{n,S}^{\Delta}-u_{n,T}^{\Delta}}_2$.
Multiplier minoration and \eqref{eq:studentized-packing-event}, intersected
with $\hat\sigma_0\le C\sigma_0$, give the same lower bound after
multiplication by the common positive factor $\sigma_0/\hat\sigma_0$.
This proves \eqref{eq:studentized-packing-lower}.  The least-squares statement
follows from \eqref{eq:least-squares-transfer}.
\end{proof}

Put
\begin{equation*}
        g_n=\frac{\mathbf X^\top\mathbf y}{\sqrt n},
        \qquad
        \hat q_{n,S}=\iota_S\hat C_{SS}^{\dagger}g_{n,S},
        \qquad \hat q_{n,\emptyset}=0.
\end{equation*}
Under the randomized null, $g_n=Z_{2n}$.

\begin{lemma}[Score-to-direction identity]
\label{lem:studentized-score-direction}
For supports $S,T$ with nonzero outcome residuals,
\begin{align}
        (I-M_S)\mathbf y&=\frac{\mathbf X\hat q_{n,S}}{\sqrt n}, \notag\\
        \norm{M_S\mathbf y}_2^2
        &=\norm{\mathbf y}_2^2-\norm{\hat q_{n,S}}_{\hat C}^2,
        \label{eq:score-projection-identity}\\
 \norm{u_{n,S}^{\Delta}-u_{n,T}^{\Delta}}_2^2
 &= (\nu_{n,S}-\nu_{n,T})^2 \notag\\
 &\quad+\frac{\nu_{n,S}\nu_{n,T}}
 {\norm{M_S\mathbf y}_2\norm{M_T\mathbf y}_2}
 \left[
 \norm{\hat q_{n,S}-\hat q_{n,T}}_{\hat C}^2
 -\{\norm{M_S\mathbf y}_2-\norm{M_T\mathbf y}_2\}^2
 \right].
 \label{eq:studentized-score-direction-bridge}
\end{align}
\end{lemma}

\begin{proof}
The Moore--Penrose identity gives
\begin{equation*}
        \frac{\mathbf X\hat q_{n,S}}{\sqrt n}
        =\mathbf X_S(\mathbf X_S^\top\mathbf X_S)^\dagger
        \mathbf X_S^\top\mathbf y=(I-M_S)\mathbf y.
\end{equation*}
The first display follows from orthogonal projection.  Moreover,
\begin{equation*}
        \norm{M_S\mathbf y-M_T\mathbf y}_2
        =\norm{\hat q_{n,S}-\hat q_{n,T}}_{\hat C}.
\end{equation*}
For nonzero vectors $a,b$,
\begin{equation*}
 \left\|\frac a{\norm a}-\frac b{\norm b}\right\|^2
 =\frac{\norm{a-b}^2-(\norm a-\norm b)^2}
 {\norm a\,\norm b}.
\end{equation*}
Apply this identity to $a=M_S\mathbf y$ and $b=M_T\mathbf y$, then use
\begin{equation*}
 \norm{\nu_{n,S}u_{n,S}-\nu_{n,T}u_{n,T}}_2^2
 =(\nu_{n,S}-\nu_{n,T})^2
 +\nu_{n,S}\nu_{n,T}\norm{u_{n,S}-u_{n,T}}_2^2.
\end{equation*}
\end{proof}

\begin{lemma}[Score packing implies direction packing]
\label{lem:score-direction-packing}
Let $\mathcal P$ be a family of size-$s$ supports and put
$r_{n,S}=\norm{M_S\mathbf y}_2$.  Suppose, for fixed constants
$c_r,c_d,\kappa>0$ and $K_q<\infty$, that
\begin{align}
 \min_{S\in\mathcal P\cup\{\emptyset\}}r_{n,S}
 &\ge c_r\norm{\mathbf y}_2,
 \label{eq:studentized-residual-floor}\\
 \min_{\substack{S,T\in\mathcal P\\S\ne T}}
 \norm{\hat q_{n,S}-\hat q_{n,T}}_{\hat C}
 &\ge d_n,
 \qquad
 \max_{S\in\mathcal P}\norm{\hat q_{n,S}}_{\hat C}
 \le K_qd_n,
 \label{eq:studentized-score-geometry}\\
 (r_{n,S}-r_{n,T})^2
 &\le(1-\kappa)
 \norm{\hat q_{n,S}-\hat q_{n,T}}_{\hat C}^2,
 \quad S\ne T,\quad S,T\in\mathcal P,
 \label{eq:studentized-angular-condition}\\
 \frac{d_n}{\norm{\mathbf y}_2}
 &\ge c_d\frac{s+1}{n}.
 \label{eq:studentized-rank-scale}
\end{align}
If $\delta\le\delta_0$ for a sufficiently small fixed constant, then
\begin{equation}\label{eq:studentized-direction-packing}
 \min_{\substack{S,T\in\mathcal P\\S\ne T}}
 \norm{u_{n,S}^{\Delta}-u_{n,T}^{\Delta}}_2
 \ge c\frac{d_n}{\norm{\mathbf y}_2},
 \qquad
 \max_{S\in\mathcal P}\norm{u_{n,S}^{\Delta}}_2
 \le C\frac{d_n}{\norm{\mathbf y}_2}.
\end{equation}
\end{lemma}

\begin{proof}
For $|S|\le s$, small $\delta$ gives
\begin{equation*}
 0<c\le\nu_{n,S}\le C,
 \qquad
 |\nu_{n,S}-\nu_{n,T}|\le C\frac{s+1}{n}.
\end{equation*}
Equations~\eqref{eq:studentized-score-direction-bridge} and
\eqref{eq:studentized-angular-condition}, together with
$r_{n,S}r_{n,T}\le\norm{\mathbf y}_2^2$, give the separation in
\eqref{eq:studentized-direction-packing}.  Taking $T=\emptyset$ in
\eqref{eq:studentized-score-direction-bridge}, dropping the nonpositive
radial term in its bracket, and using
\eqref{eq:studentized-residual-floor} gives
\begin{equation*}
 \norm{u_{n,S}^{\Delta}}_2
 \le C\left\{\frac{s+1}{n}
 +\frac{\norm{\hat q_{n,S}}_{\hat C}}{\norm{\mathbf y}_2}\right\}.
\end{equation*}
Equations~\eqref{eq:studentized-score-geometry} and
\eqref{eq:studentized-rank-scale} prove the radius bound.
\end{proof}

The angular condition has the equivalent residual-energy form
\begin{equation}\label{eq:studentized-energy-balance}
 \left|\norm{\hat q_{n,S}}_{\hat C}^2
 -\norm{\hat q_{n,T}}_{\hat C}^2\right|
 \le\sqrt{1-\kappa}\,(r_{n,S}+r_{n,T})
 \norm{\hat q_{n,S}-\hat q_{n,T}}_{\hat C}.
\end{equation}
Indeed,
$r_{n,S}^2=\norm{\mathbf y}_2^2-\norm{\hat q_{n,S}}_{\hat C}^2$.
A simpler sufficient condition is
\begin{equation}\label{eq:studentized-simple-energy-balance}
 \max_{S\in\mathcal P}\norm{\hat q_{n,S}}_{\hat C}
 \le\sqrt{1-\kappa}
 \min_{S\in\mathcal P}r_{n,S}.
\end{equation}
This follows from the triangle inequality because the left side of
\eqref{eq:studentized-energy-balance} is at most
\begin{equation*}
 2\max_{R\in\mathcal P}\norm{\hat q_{n,R}}_{\hat C}
 \norm{\hat q_{n,S}-\hat q_{n,T}}_{\hat C},
\end{equation*}
whereas $r_{n,S}+r_{n,T}$ is at least twice the minimum residual norm.

Under the randomized null, $g_n=Z_{2n}$.  Relative restricted isometry
and inverse perturbation on every
$2s$-coordinate union gives, uniformly over the pool in
Lemma~\ref{lem:top-s-score-pool},
\begin{equation*}
 \norm{\hat q_{n,S}-q_{n,S}}_C
 \le C\delta\norm{q_{n,S}}_C,
 \qquad
 \norm z_{\hat C}=\{1+O(\delta)\}\norm z_C
 \quad\text{for }|\operatorname{supp}(z)|\le2s.
\end{equation*}
Thus local isotropy and $\delta+\rho_s(C)=o(1)$ transfer
\eqref{eq:top-s-score-pool} to the empirical geometry in
\eqref{eq:studentized-score-geometry}.  The outcome floor in
\eqref{eq:t-floors} and sparse concentration also give
\eqref{eq:studentized-residual-floor} outside an event of probability
$Ce^{-cm}$.  On
$\norm{\mathbf y}_2\asymp\sqrt{n\Var(y)}$, a noise pool with
$b_n\asymp\sqrt{\Var(y)L}$ has $d_n\asymp\sqrt{\Var(y)m}$ and therefore
gives $\lambda_n\asymp\delta$.  Since
$(s+1)/n=o(\delta)$ when $m=o(n)$, it satisfies
\eqref{eq:studentized-rank-scale}; its score radius is $o(r_{n,S})$, so
\eqref{eq:studentized-simple-energy-balance} holds.  A signal pool with
$b_n\asymp\sqrt n\,\Psi_s/\sqrt s$ has
$d_n\asymp\sqrt n\,\Psi_s$ and gives
$\lambda_n\asymp\sqrt{\theta_s}$ provided
\eqref{eq:studentized-angular-condition} holds and
$\sqrt{\theta_s}\gtrsim(s+1)/n$.  The latter rank-scale condition is
automatic in the signal-dominated regime
$\sqrt{\theta_s}\gtrsim\delta$.  The signal pool and its
energy balance are multiplicity conditions, not consequences of the
signal norm alone.

\end{document}
